\chapter{Upper bounds: P-LCP, UEOPL and a subexponential algorithm}\label{ch:upper}

\paragraph{Idea.} This chapter gives three upper bounds, by three routes.
\begin{itemize}[nosep]
\item \emph{Through the market} (Section~\ref{sec:plcp}). The equilibrium conditions of (M1) are a
linear complementarity problem whose matrix is almost a P-matrix. Adding a tiny multiple of the identity
makes it one, and the support of the perturbed solution fixes an LP whose every feasible point is an
exact equilibrium. So (M1),
and with it ARRIVAL, reduce many-one to P-LCP.
\item \emph{Directly from the game to one P-matrix LCP} (Section~\ref{sec:direct-lcp}). The BNE conditions
of the game itself, after a row shift, form one LCP whose matrix is $P_0$; a tiny diagonal perturbation
makes it a P-matrix. This is a many-one reduction of exact BNE to P-LCP, and it puts the game in
$\mathsf{UEOPL}$.
\item \emph{Through one contraction} (Sections~\ref{sec:contraction}--\ref{sec:subexp}). The support
and coverage labels for an equilibrium LP are read from the unique fixed point of an explicit piecewise-linear map
$F:[0,1]^{n+1}\to[0,1]^{n+1}$ that contracts the $\ell_1$ norm. The map is built once, in
Section~\ref{sec:contraction}, and used twice: read as a PL-Contraction instance it puts exact BNE in
$\mathsf{UEOPL}$ (Section~\ref{sec:ueopl}); handed to the block search of Ghorbani, Lopes and Mnich it
gives an $L^{O(\sqrt n)}$ algorithm (Section~\ref{sec:subexp}).
\end{itemize}

\section{The market reduces to P-matrix LCP}\label{sec:plcp}

\paragraph{Idea.} Write the four conditions (a)--(d) of Definition~\ref{def:m1} as complementarity
between each variable and its slack. The resulting matrix $M$ has a zero diagonal, so it is not a
P-matrix, but all its principal minors are nonnegative because the market's bipartite graph turns
every nonzero minor into a single matching term. Then $M+\varepsilon I$ is a P-matrix. If $\varepsilon$
is small enough, a Cramer-rule gap shows that the support of the unique perturbed solution fixes an LP that is
feasible and whose every feasible point is an exact equilibrium.

The scheme is that of Adler and Verma~\cite[\S4]{av11}. Their Theorem~4.3 covers nondegenerate LCPs whose matrix lies in a Lemke-processable class of semimonotone matrices: for every sufficiently small $\varepsilon>0$, a solution of $\mathrm{LCP}(q,M+\varepsilon I)$ yields in polynomial time, by linear programs on its support, a solution of $\mathrm{LCP}(q,M)$ or a certificate that none exists. They note that $P_0\cap Q_0$-LCPs thereby reduce to P-LCP, and every (M1) matrix lies in $P_0\cap Q_0$ (Proposition~\ref{plcp:prop:q0}). The proof below is self-contained. It adds an explicit $\varepsilon$ of polynomial bit length, which a polynomial-time many-one reduction needs, and it assumes no nondegeneracy.

\paragraph{Example.} For one buyer and one good the LCP has three coordinates $(f,u,P)$, and
Example~\ref{plcp:ex:negative} below works it out. It also shows that $\varepsilon$ cannot be arbitrary: with
$\varepsilon=1$ the perturbed solution has the wrong support and the decoding LP is infeasible, while for every $\varepsilon<1$ the support is right and the LP returns the equilibrium. In this section all (M1) slopes need only be positive, $b_{ij}>0$.

\subsection{The LCP formulation and its matrix}

Fix an (M1) market with edge set $E$. Index coordinates by edges, buyers and goods:
$N=|E|+|I|+|J|$, and $z=(f,u,P)\in\mathbb R^E\times\mathbb R^I\times\mathbb R^J$. Define $w=q+Mz$ by
\[
  w_{ij}=u_i+b_{ij}P_j-a_{ij}\ \ ((i,j)\in E),\qquad
  w_i=\mu_i-\sum_{j}f_{ij}\ \ (i\in I),\qquad
  w_j=c_j-\sum_{i}f_{ij}\ \ (j\in J).
\]
So $q=(-a,\mu,c)$ and, with the vertex set $V=I\cup J$,
\[
  M=\begin{pmatrix}0&B\\-K^{\top}&0\end{pmatrix},
\]
where $B,K\in\mathbb Q^{E\times V}$: the row of edge $(i,j)$ has entries $1$ (column $i$) and $b_{ij}$
(column $j$) in $B$, entries $1$ in both columns $i$ and $j$ in $K$, and zeros elsewhere. The pairs
$(f_{ij},w_{ij})$, $(u_i,w_i)$ and $(P_j,w_j)$ are complementary.

\begin{lemma}[LCP formulation]\label{plcp:lem:formulation}
$(f,u,P)\in\mathrm{SOL}(q,M)$ if and only if $(P,v,f)$ with $v=u$ is an (M1) equilibrium.
\end{lemma}

\begin{proof}
$z\ge0$ is $f,v,P\ge0$. $w_{ij}\ge0$ is (a), and $f_{ij}w_{ij}=0$ is (b). $w_i\ge0$ with $u_iw_i=0$ is
(c), and $w_j\ge0$ with $P_jw_j=0$ is (d).
\end{proof}

Both signs of $a_{ij}$ are allowed, and the edge set $E$ is arbitrary. The diagonal of $M$ is zero, so
$M$ is not a P-matrix.

\begin{lemma}[$P_0$]\label{plcp:lem:p0}
For every (M1) market (any edge set $E$ and any positive slopes), every principal minor of $M$ is
nonnegative.
\end{lemma}

\paragraph{Proof roadmap.} Three cases for a principal submatrix with edge rows $E'$ and vertex columns
$V'$: (1) if $|E'|\ne|V'|$ the rank is too small; (2) if $E'$ contains a cycle, the alternating sum of
its rows of $K$ vanishes; (3) if $E'$ is a forest, both determinants have the same single matching term,
and the minor is a product of positive entries.

\begin{proof}
A principal submatrix of $M$ is determined by a set $E'\subseteq E$ of edge coordinates and a set
$V'\subseteq V$ of vertex coordinates, and has the form
$\left(\begin{smallmatrix}0&X\\Y&0\end{smallmatrix}\right)$ with $X=B_{E'V'}$ and $Y=-(K_{E'V'})^{\top}$.
If $|E'|\ne|V'|$, its rank is at most $2\min(|E'|,|V'|)<|E'|+|V'|$, so the minor is $0$. If
$|E'|=|V'|=k$, the minor equals $(-1)^{k}\det X\det Y=(-1)^k\det B_{E'V'}\cdot(-1)^k\det K_{E'V'}
=\det B_{E'V'}\det K_{E'V'}$.

If the edges $E'$ contain a cycle of the bipartite graph $(V,E)$, sum the rows of $K$ indexed by the
edges of the cycle with alternating signs $\pm1$. The cycle has even length, every vertex on it is an
endpoint of exactly two consecutive cycle edges, which receive opposite signs, and vertices off the cycle
receive nothing. So the combination is $0$, also after deleting the columns outside $V'$; hence
$\det K_{E'V'}=0$ and the minor is $0$.

If $E'$ is a forest, consider the bipartite graph $G'$ whose two sides are the rows $E'$ and the
columns $V'$, with $e$ adjacent to $v$ when $v$ is an endpoint of $e$. It is a subgraph of the
subdivision of a forest, hence a forest, and therefore has at most one perfect matching (the symmetric
difference of two distinct perfect matchings contains a cycle). The nonzero entries of $B_{E'V'}$ and of
$K_{E'V'}$ occupy exactly the positions of the edges of $G'$, because $b_{ij}>0$. A nonzero term in the
Leibniz expansion of either determinant is a perfect matching of $G'$. So either both determinants are
$0$, or both consist of the same single term $\sigma$: $\det K_{E'V'}=\operatorname{sgn}\sigma$ and
$\det B_{E'V'}=\operatorname{sgn}\sigma\cdot\beta$, where $\beta$ is a product of entries equal to $1$
or to some $b_{ij}>0$. The minor is then $\beta>0$.
\end{proof}

\begin{lemma}[Diagonal perturbation]\label{plcp:lem:pmatrix}
For every $P_0$-matrix $M$ and every $\varepsilon>0$, $M+\varepsilon I$ is a P-matrix. In particular
this holds for the (M1) matrix.
\end{lemma}

\begin{proof}
By multilinearity of the determinant in the columns,
\[\det(M_{SS}+\varepsilon I)=\sum_{T\subseteq S}\varepsilon^{|S\setminus T|}\det M_{TT}.\] Every term is
nonnegative, and the term $T=\emptyset$ equals $\varepsilon^{|S|}>0$.
\end{proof}

\subsection{The reduction}

In this section only, define for a given (M1) market
\[
  H=\max\Bigl(0,\max_{(i,j)\in E}a_{ij}\Bigr),\qquad b_{\min}=\min_{(i,j)\in E}b_{ij},\qquad
  \bar Z=1+\max\Bigl\{\max_i\mu_i+H,\ \frac{H}{b_{\min}}\Bigr\},
\]
$d=N+1$, $D$ the least common multiple of all denominators of entries of $M$ and $q$, and
$C=D\cdot\max(1,\max_{k,l}|M_{kl}|,\max_k|q_k|)$, an integer.

\begin{lemma}[Uniform bound]\label{plcp:lem:bound}
For every $0<\varepsilon\le1$, every solution $z^\varepsilon=(f,u,P)$ of $\mathrm{LCP}(q,M+\varepsilon I)$
satisfies $0\le z^\varepsilon_k\le\bar Z-1$ for all $k$.
\end{lemma}

\begin{proof}
Write the perturbed slacks $w^\varepsilon=q+(M+\varepsilon I)z^\varepsilon$. If $u_i>0$, then
$w^\varepsilon_i=0$, i.e.\ $\sum_jf_{ij}=\mu_i+\varepsilon u_i>0$, so some $f_{ij}>0$, and then
$w^\varepsilon_{ij}=0$ reads $u_i+b_{ij}P_j+\varepsilon f_{ij}=a_{ij}$; since all terms on the left are
nonnegative, $u_i\le a_{ij}\le H$. If $P_j>0$, then $\sum_if_{ij}=c_j+\varepsilon P_j>0$, so some
$f_{ij}>0$, and the same equation gives $b_{ij}P_j\le a_{ij}$, so $P_j\le H/b_{\min}$. Finally
$w^\varepsilon_i\ge0$ gives $f_{ij}\le\sum_{j'}f_{ij'}\le\mu_i+\varepsilon u_i\le\mu_i+H$.
\end{proof}

For $A\subseteq[N]$ write $\bar A=[N]\setminus A$, and consider the LP
\[
  \mathrm{LP}(A):\quad\min\ \rho\quad\text{s.t.}\quad z\ge0,\ \rho\ge0,\ z_{\bar A}=0,\
  q+Mz\ge-\rho\mathbf 1,\ (q+Mz)_A\le\rho\mathbf 1 .
\]

\begin{lemma}[Uniform gap]\label{plcp:lem:gap}
For every $A\subseteq[N]$, $\mathrm{LP}(A)$ has an optimum $\rho^*_A$, and either $\rho^*_A=0$ or
$\rho^*_A\ge(dC)^{-d}$.
\end{lemma}

\begin{proof}
$z=0$ with $\rho$ large is feasible, and $\rho\ge0$ bounds the objective, so an optimum exists. The
feasible region lies in the orthant $\{z\ge0,\rho\ge0\}$, so it is a pointed polyhedron and the optimum
is attained at a vertex. Multiply every constraint by $D$: all coefficients and right-hand sides become
integers of absolute value at most $C$ (the coefficient of $\rho$ becomes $D\le C$). A vertex is the
unique solution of $d$ linearly independent tight constraints, so by Cramer's rule
$\rho^*_A=\Delta'/\Delta$ with integers $\Delta',\Delta$ and $0<|\Delta|\le d!\,C^d\le(dC)^d$. If
$\rho^*_A>0$ then $|\Delta'|\ge1$ and $\rho^*_A\ge(dC)^{-d}$.
\end{proof}

\begin{theorem}[(M1) $\le_m$ P-LCP]\label{plcp:thm:main}
\statusrefereed
Let an (M1) market be given, with $q,M,\bar Z,d,C$ as above, and put
\[
  \varepsilon=\frac{1}{2\lceil\bar Z\rceil(dC)^d}.
\]
Then $M+\varepsilon I$ is a P-matrix. Let $z^\varepsilon$ be the unique solution of
$\mathrm{LCP}(q,M+\varepsilon I)$ and $A=\{k:z^\varepsilon_k>0\}$. Then the system
\[
  z\ge0,\qquad z_{\bar A}=0,\qquad q+Mz\ge0,\qquad(q+Mz)_A=0 \tag{$\ast$}
\]
is feasible, and every feasible point $z=(f,u,P)$ of $(\ast)$ gives an (M1) equilibrium $(P,u,f)$.
Moreover $\log_2(1/\varepsilon)=O(L^2)$, the instance $(q,M+\varepsilon I)$ has $O(L^3)$ bits and is
computable in polynomial time, and $(\ast)$ is a rational LP feasibility problem of polynomial size.
Hence the (M1) equilibrium problem has a polynomial-time many-one reduction to P-LCP, whose solution map
is one rational LP.
\end{theorem}

\begin{proof}
The proof has three parts: feasibility of $(\ast)$ via the gap lemma, soundness via the LCP
formulation, and size.

$M+\varepsilon I$ is a P-matrix by Lemmas~\ref{plcp:lem:p0} and~\ref{plcp:lem:pmatrix}, so
$z^\varepsilon$ exists and is unique (Proposition~\ref{prop:classical}(i)).

\emph{Feasibility of $(\ast)$.} Since $\varepsilon\le1$, Lemma~\ref{plcp:lem:bound} gives
$z^\varepsilon\le(\bar Z-1)\mathbf 1$. As $q+(M+\varepsilon I)z^\varepsilon\ge0$, every original slack
satisfies $(q+Mz^\varepsilon)_k\ge-\varepsilon z^\varepsilon_k\ge-\varepsilon\bar Z$. For $k\in A$,
complementarity in the perturbed LCP gives $(q+Mz^\varepsilon)_k=-\varepsilon z^\varepsilon_k\le0$. And
$z^\varepsilon_{\bar A}=0$ by definition of $A$. Hence $(z^\varepsilon,\rho)$ with
$\rho=\varepsilon\bar Z$ is feasible for $\mathrm{LP}(A)$, and
\[
  \varepsilon\bar Z=\frac{\bar Z}{2\lceil\bar Z\rceil(dC)^d}\le\frac{1}{2(dC)^d}<(dC)^{-d}.
\]
By Lemma~\ref{plcp:lem:gap}, $\rho^*_A=0$; an optimal $z$ of $\mathrm{LP}(A)$ is feasible for $(\ast)$.

\emph{Soundness.} A feasible $z$ of $(\ast)$ has $z\ge0$ and $w=q+Mz\ge0$, with $w_k=0$ for $k\in A$ and
$z_k=0$ for $k\notin A$. So $z\in\mathrm{SOL}(q,M)$, and Lemma~\ref{plcp:lem:formulation} applies. This
holds even if some coordinates in $A$ vanish at $z$, so the decoder never needs to identify a particular
equilibrium.

\emph{Size.} $N=|E|+|I|+|J|=O(L)$, so $d=O(L)$. $D$ divides the product of all input denominators, so
$\log D=O(L)$, and $\log C=O(L)$; also $\log\lceil\bar Z\rceil=O(L)$. Hence
$\log_2(1/\varepsilon)=1+\log_2\lceil\bar Z\rceil+d\log_2(dC)=O(L^2)$, and the $N^2$ entries of
$M+\varepsilon I$ together with $q$ take $O(N^2+NL^2)=O(L^3)$ bits. All of $H$, $b_{\min}$, $\bar Z$,
$D$, $C$ and $\varepsilon$ are computed directly from the input. The system $(\ast)$ has the original
(polynomial-size) data, and a rational solution of polynomial size is found in polynomial time by the
ellipsoid or an interior-point method.
\end{proof}

The proof uses only that $z^\varepsilon$ is \emph{some} solution of the perturbed LCP. Since a P-matrix
LCP always has a solution, it also shows that every (M1) market (with $b_{ij}>0$) has an equilibrium. The
perturbation $\varepsilon$ is exponentially small (polynomially many bits), and the decoding uses only
the support of the perturbed solution, not its values.

\begin{example}[The size of $\varepsilon$ matters]\label{plcp:ex:negative}
Take one buyer and one good with $\mu=1$, $c=2$, $a=b=1$. The unique equilibrium is $f=1$, $v=1$,
$P=0$. At $\varepsilon=1$ the perturbed LCP has the solution $(f,u,P)=(1,0,0)$: indeed
$w_f=u+P-1+f=0$, $w_u=1-f+u=0$ and $w_P=2-f+P=1$. Its support $A=\{f\}$ forces $u=P=0$ in $(\ast)$, and
then $w_f=-1<0$, so $(\ast)$ is infeasible. For $0<\varepsilon<1$ the perturbed solution is
$u=(1-\varepsilon)/(1+\varepsilon^2)>0$, $f=1+\varepsilon u$, $P=0$, with support $\{f,u\}$, and $(\ast)$
returns the equilibrium.
\end{example}

\begin{corollary}[(M1) is in UEOPL]\label{plcp:cor:ueopl}
\statusrefereed
The (M1) equilibrium problem lies in $\mathsf{UEOPL}$.
\end{corollary}

\begin{proof}
The instance $(q,M+\varepsilon I)$ produced by Theorem~\ref{plcp:thm:main} is a P-matrix instance, so it
has no PV2 violation. Compose with the promise-preserving reduction of~\cite[Thm.~81]{fgms18} from P-LCP
with PV2 violations to $\mathsf{UEOPL}$, and decode by $(\ast)$; $\mathsf{UEOPL}$ is closed under
many-one reductions.
\end{proof}

\subsection{Consequences for ARRIVAL and for the game}

\begin{corollary}[ARRIVAL $\le_m$ P-LCP]\label{plcp:cor:arrival}
\statusrefereed
The ARRIVAL decision problem has a polynomial-time many-one reduction to P-LCP. The P-LCP instance
produced is a P-matrix instance of $\mathrm{poly}(n)$ bits. Because $M+\varepsilon I$ is a P-matrix, this is a reduction to the
promise version of P-LCP.
\end{corollary}

\begin{proof}
Map the ARRIVAL instance to the market of Theorem~\ref{thm:arrival-m1}, which has $\mathrm{poly}(n)$
size and positive slopes, and that market to $(q,M+\varepsilon I)$ by Theorem~\ref{plcp:thm:main}. Given
the solution $z^\varepsilon$, solve $(\ast)$ to get an exact (M1) equilibrium, and answer \yes{} if and
only if the non-terminal buyers' load on $T_{\yes}$ is $1$. By Theorem~\ref{thm:arrival-m1} this is
correct, because its conclusion holds for \emph{every} equilibrium. There is one instance map and one
solution map (an LP plus a readout), so the composite is many-one.
\end{proof}

\paragraph{The search version.}\label{par:search-arrival}
The same P-LCP answer gives the switching-flow certificate required by
\textsc{S-Arrival} in \cite[Definition~13]{ghhkms18}. Form the specified
graph $G'$ by adding $\bar d$, redirecting both exits of every original
vertex with no path to $d$ to $\bar d$, and making $d,\bar d$ absorbing.
Use these two vertices as the terminals in the market reduction. Every
remaining switch can reach a terminal. If $o=d$, the zero flow solves the
input and the decoder can ignore a fixed trivial P-LCP query.
Otherwise read the incoming masses $x_v$ of Theorem~\ref{thm:compose}
from the recovered market equilibrium. On each labeled exit $(v,i)$
output $h_i(\lfloor x_v\rfloor)$, setting the terminal loops to zero.
Parts (i)--(ii) of the proof of Lemma~\ref{lem:sh-decoder} show that flooring
preserves the switch balances and both terminal inflows. The two
labeled counts at each switch are $\lceil\lfloor x_v\rfloor/2\rceil$
and $\lfloor\lfloor x_v\rfloor/2\rfloor$. If the output stores one
edge per ordered vertex pair, add the two counts when the successors
coincide. Conservation is unchanged and parity on a repeated edge
holds automatically. These are the conservation and parity requirements for a
switching flow to the reached terminal. All counts have $O(n)$ bits
by Theorem~\ref{thm:compose}, and flooring and splitting take polynomial
time. The certificate may contain extra circulations. Definition~14
of the same paper asks instead for a run profile.

G\"artner et al.~\cite[Section~6]{ghhkms18} asked about a reduction
between ARRIVAL and P-LCP. Corollary~\ref{plcp:cor:arrival} supplies
the decision direction, and the preceding decoder supplies the
switching-flow search direction. Earlier reductions place the search
problem in $\mathsf{UEOPL}$~\cite{ghhkms18,fgms18}. P-LCP also lies in
that class~\cite{fgms18}. The construction through the market gives
the instance map and witness map relating the two problems. Composing
with the equivalence of P-LCP and P-Lin-Bellman
from~\cite{bfgssw24} also gives a search reduction to P-Lin-Bellman.

\begin{remark}[The market route]\label{plcp:cor:bne}
By Proposition~\ref{prop:mk-g2m} and Theorem~\ref{plcp:thm:main}, exact BNE also has a polynomial-time Turing reduction to P-LCP through the market; the direct reduction of Section~\ref{sec:direct-lcp} is stronger.
\end{remark}

\begin{proof}
Run the oracle algorithm of Proposition~\ref{prop:mk-g2m} (Theorem~\ref{thm:mk-main}(b)), whose every
call is an ordinary (M1) market with positive masses, capacities and slopes, and answer each (M1) call by
one P-LCP call and one LP (Theorem~\ref{plcp:thm:main}). Each call has polynomial size, and there are
polynomially many.
\end{proof}

The composition is Turing, not many-one, because the game-to-market reduction is adaptive. In
particular, composing the many-one reduction from ARRIVAL to the game with
Corollary~\ref{plcp:cor:bne} would give only a Turing reduction from ARRIVAL to P-LCP;
Corollary~\ref{plcp:cor:arrival} obtains a many-one reduction by bypassing the game and going through
(M1) directly. Corollary~\ref{plcp:cor:bne} by itself does not place the game in $\mathsf{UEOPL}$; that needs
a many-one reduction. Section~\ref{sec:direct-lcp} gives one directly, and Section~\ref{sec:ueopl} gives a second,
independent route.

\subsection{Lemke's algorithm on the unperturbed matrix}

One may ask whether the perturbation is needed at all. $M$ is not a P-matrix, and $\mathrm{LCP}(0,M)$
has nonzero solutions (for one buyer and one good with $a=\mu=c=0$, $u>0$ and $f=P=0$ is one), so the
classical finite-termination criteria for Lemke's algorithm that exclude such solutions do not apply. A
different criterion does.

\begin{proposition}[$Q_0$]\label{plcp:prop:q0}
\statusrefereed
For every (M1) matrix $M$ (any edge set, positive slopes), $M\in P_0\cap Q_0$.
\end{proposition}

\paragraph{Proof roadmap.} (1) Feasibility of $q$ forces $\mu,c\ge0$; (2) delete zero-mass buyers and
zero-capacity goods; (3) Kakutani's theorem on a price box gives an equilibrium of the rest; (4) restore
the deleted buyers and goods with zero flow and suitable prices and utilities.

\begin{proof}
$P_0$ is Lemma~\ref{plcp:lem:p0}. Let $q=(-a,\mu,c)$ be feasible, i.e.\ some $z\ge0$ has $q+Mz\ge0$.
The buyer and good rows of $q+Mz$ are $\mu_i-\sum_jf_{ij}$ and $c_j-\sum_if_{ij}$, so feasibility forces
$\mu\ge0$ and $c\ge0$, while the intercepts $a$ are arbitrary. We show that such a market, possibly with
zero masses or capacities, has an equilibrium, which by Lemma~\ref{plcp:lem:formulation} solves
$\mathrm{LCP}(q,M)$.

Delete buyers with $\mu_i=0$ and goods with $c_j=0$. For the remaining market take a price box
$[0,\bar B]^J$ with $\bar B>\max(0,\max_{ij}a_{ij}/b_{ij})$. At prices $P$, let each buyer distribute its
mass over its best replies, the outside option included; at an allocation $f$, let the prices be chosen to
maximize $\sum_jP_j(\sum_if_{ij}-c_j)$ over the box. Both correspondences have nonempty convex values and
closed graphs, so Kakutani's theorem on the product of the allocation simplices and the price box gives a
joint fixed point. No good has price $\bar B$ there: its payoff is negative for every buyer, so it has
zero load against a positive capacity, and the unique price best reply is $0$. Hence the fixed point
satisfies (a)--(d). (If one side or every edge was deleted, zero flows with zero prices and utilities
solve the remaining market; a maximum over an empty family is $0$ here.) Restore deleted goods with zero
flow and price at least $\max(0,\max_ia_{ij}/b_{ij})$ (price zero for a good that no buyer lists),
and deleted buyers with zero flow and utility $\max(0,\max_j(a_{ij}-b_{ij}P_j))$. All slacks and
complementarity conditions hold.
\end{proof}

For $P_0$ matrices, Aganagi\'c and Cottle~\cite{ac87} give a constructive characterization of $Q_0$.
What matters here is the following standard fact: Lemke's algorithm with lexicographic pivoting always
terminates, either at a solution or on a secondary ray, and for $M\in P_0\cap Q_0$ a secondary ray
certifies that $q$ is infeasible (see~\cite[\S2, Remarks~2 and~5]{av13}). Since a market with positive
masses and capacities gives a feasible $q$ (take $f=0$, $P=0$ and $u_i=H$ for every buyer), Lemke's algorithm applied to $\mathrm{LCP}(q,M)$
terminates at an (M1) equilibrium. No polynomial bound on the number of pivots is known, neither on $M$
nor on $M+\varepsilon I$.

\begin{remark}[No secondary ray, directly; a sketch]\label{plcp:rem:noray}
For positive masses and capacities there is also a direct argument. For a strictly positive rational
covering vector $h$, consider $w=q+Mz+\tau h$, $\tau\ge0$, with all pairs complementary. At fixed $\tau$
this is the (M1) market with intercepts $a_{ij}-\tau h_{ij}$, masses $\mu_i+\tau h_i$ and capacities
$c_j+\tau h_j$. For $\tau>T=\max(0,\max_{ij}a_{ij}/h_{ij})$ all intercepts are negative, so $f=0$, and the
positive margins force $u=P=0$. For $\tau\le T$ the flows are bounded by the row inequalities, and a
positive utility or price needs a used edge, which gives the bounds of Lemma~\ref{plcp:lem:bound}. So the
only unbounded ray is the primary ray $z=0$, $\tau\ge T$. Under the lexicographic rule the primary path
then ends at $\tau=0$; degenerate inputs are handled by a symbolic perturbation of $q$ and a limiting
argument over finitely many bases.
\end{remark}

\section{The game reduces directly to one P-matrix LCP}\label{sec:direct-lcp}

This section proves the main upper bound for the game: exact BNE search has a polynomial-time many-one reduction to one P-LCP instance whose matrix is a genuine P-matrix (Theorem~\ref{thm:direct-lcp}). Unlike Section~\ref{sec:plcp}, it does not go through the market and its adaptive reduction.

\paragraph{Idea.} Write the BNE conditions of the game itself as one linear complementarity problem,
with one variable per type--target pair, per type, per target (two) and one for the budget. A forest and
cofactor argument on two incidence matrices shows that its matrix $M$ is $P_0$. Then every $M+\delta I$
with $\delta>0$ is a P-matrix, and a uniform bound together with a Cramer-rule gap shows that one
explicit $\delta$ of polynomial bit length suffices: the support of the unique solution of the
perturbed LCP determines a linear program whose solutions are exact equilibria. The reduction allows zero
slopes, zero defender gains, payoff ties and every budget.

\paragraph{Notation.} The game is that of Definition~\ref{def:game}, in the coordinates of
Section~\ref{sec:pre-game}: types $i\in[m]$ with priors $p_i$, targets $j\in[n]$, budget $r$, intercepts
$a_{ij}$, slopes $b_{ij}\ge0$ and defender gains $d_j\ge0$ as in~\eqref{eq:dictionary}, coverage $x$,
attack masses $f_{ij}$, loads $w_j=\sum_if_{ij}$ and scores $\pi_j=d_jw_j$. LCPs, P- and $P_0$-matrices
and principal submatrices $M_{SS}$ are as in Section~\ref{sec:pre-classes}, and many-one reductions as in
Section~\ref{sec:pre-red}. For $\mathrm{LCP}(q,M)$ the coordinate $(q+Mz)_k$ is the \emph{slack} paired
with $z_k$. The symbols $M$, $q$, $N$, $B$, $D$ and the constants $U$, $W$, $D_*$, $S$, $K$, $C$, $Q$ are
local to this section and differ from Section~\ref{sec:plcp}; in particular $U$ without indices is a constant,
not the payoff $U_{ij}(x)$, and the LCP variable $\lambda$ is the cutoff of Lemma~\ref{lem:cutoff} plus
one. Appendix~\ref{app:sources} records its history, and Appendix~\ref{app:status} its checks.

\subsection{BNE in marginal coordinates}\label{app-dl:sec:bne}

Let types $1,\dots,m$ be the types of positive prior, so $p_i>0$ and $\sum_ip_i=1$. Types of prior zero
are handled separately (Section~\ref{app-dl:sec:decode}).

\begin{definition}[BNE in marginal coordinates]\label{app-dl:def:bne}
A pair $(x,f)$ with $x\in[0,1]^n$, $\sum_jx_j=r$, $f\ge0$ and $\sum_jf_{ij}=p_i$ for all $i$ is a
\emph{BNE in marginal coordinates} if (a) $f_{ij}>0$ implies $a_{ij}-b_{ij}x_j=\max_k(a_{ik}-b_{ik}x_k)$,
and (b) $x$ maximizes $\sum_j\pi_jx'_j$ over all $x'\in[0,1]^n$ with $\sum_jx'_j=r$.
\end{definition}

Up to a term that does not depend on the defender's strategy, the defender's expected payoff is
$\sum_j\pi_jx_j$. Every mixed defender strategy induces a feasible $x'$, and every feasible $x'$ is
induced by an explicit mixed strategy (Lemma~\ref{app-dl:lem:decode}). Hence a defender mixed strategy
with marginals $x$, together with the attack distributions $f_{ij}/p_i$, is a BNE if and only if $(x,f)$
is a BNE in marginal coordinates. For $0<r<n$ we use the cutoff form of~(b), Lemma~\ref{lem:cutoff}(a):
condition~(b) holds if and only if there is $\lambda_0\in\mathbb R$ with
$x_j>0\Rightarrow\pi_j\ge\lambda_0$ and $x_j<1\Rightarrow\pi_j\le\lambda_0$.

\subsection{Normalization and the complementarity system}\label{app-dl:sec:lcp}

\paragraph{Edge cases.} If $r=0$, the coverage is $x=0$. If $r=n$, it is $x=(1,\dots,1)$. In both cases
the defender has a single pure strategy, and a BNE is obtained by giving each type, including every
zero-prior type, a target maximizing $a_{ij}-b_{ij}x_j$. To keep the many-one format, the instance map
sends such games to the fixed instance $M=[1]$, $q=[0]$, and the solution map ignores the answer.
\emph{From now on $0<r<n$}, and zero-prior types are dropped until Section~\ref{app-dl:sec:decode}.

\paragraph{Shifting the attacker payoffs.} For each type $i$ put
\[
  A_{ij}=a_{ij}+1+\max_k\,(b_{ik}-a_{ik}) .
\]
Then $A_{ij}-b_{ij}=1+\max_k(b_{ik}-a_{ik})-(b_{ij}-a_{ij})\ge1$, so the shifted payoff
$A_{ij}-b_{ij}x_j$ is at least $1$ on $[0,1]^n$. The shift is a constant for each type, so it preserves
every type's preferences over targets at every coverage vector.

\paragraph{Shifting the defender scores.} We add $1$ to every defender score, that is, we use the scores
$1+\pi_j$. On full-budget coverage vectors ($\sum_jx_j=r$) this changes the objective
$\sum_j\pi_jx_j$ by the constant $r$, so it does not change the set of maximizers. Its purpose is to make
all shifted scores at least $1$, which forces the budget multiplier below to be positive.

\paragraph{Variables and pairs.} The LCP has nonnegative variables $f_{ij}$ (joint attack masses),
$u_i$ (shifted best attacker utility of type $i$), $x_j$ (coverage), $t_j$ (multiplier of the cap
$x_j\le1$) and $\lambda$ (shifted defender cutoff). With $w_j=\sum_if_{ij}$, the complementary pairs are
\begin{equation}\label{app-dl:eq:lcp}
\begin{array}{c|l}
  \text{variable} & \text{paired slack}\\\hline
  f_{ij} & u_i+b_{ij}x_j-A_{ij}\\
  t_j    & 1-x_j\\
  \lambda& r-\sum_jx_j\\
  u_i    & p_i-\sum_jf_{ij}\\
  x_j    & \lambda-1+t_j-d_jw_j .
\end{array}
\end{equation}
Write this system as $\mathrm{LCP}(q,M)$. The variable order is $z=(f,t,\lambda,u,x)$, so the dimension is
\[
  N=mn+n+1+m+n=mn+m+2n+1 .
\]

\begin{lemma}[LCP formulation]\label{app-dl:lem:formulation}
\statuslean
Let $0<r<n$. (a) If $z=(f,t,\lambda,u,x)$ solves $\mathrm{LCP}(q,M)$, then $\lambda>0$,
$\sum_jx_j=r$, and $(x,f)$ is a BNE in marginal coordinates. (b) Conversely, every BNE $(x,f)$ in
marginal coordinates extends to a solution $(f,t,\lambda,u,x)$ of $\mathrm{LCP}(q,M)$.
\end{lemma}

\begin{proof}
(a) The cap slacks give $x_j\le1$, so $0\le x\le1$. The first slack gives
$u_i\ge A_{ij}-b_{ij}x_j\ge A_{ij}-b_{ij}\ge1>0$. Complementarity at $u_i$ then makes its slack zero:
$\sum_jf_{ij}=p_i$. If $f_{ij}>0$, then $u_i=A_{ij}-b_{ij}x_j$, while $u_i\ge A_{ik}-b_{ik}x_k$ for every
$k$. So positive flow uses only maximizers of the shifted, hence of the original, payoff. This is
condition~(a) of Definition~\ref{app-dl:def:bne}.

\emph{Why $\lambda>0$.} Suppose $\lambda=0$. The slack of $x_j$ is nonnegative, so
$t_j\ge1+d_jw_j\ge1>0$ for every $j$. Complementarity at $t_j$ gives $x_j=1$ for every $j$, so
$\sum_jx_j=n>r$, contradicting the nonnegative budget slack $r-\sum_jx_j$. Hence $\lambda>0$, and
complementarity at $\lambda$ gives $\sum_jx_j=r$.

\emph{Defender optimality.} Put $\lambda_0=\lambda-1$. If $x_j<1$, complementarity at $t_j$ gives
$t_j=0$, and the nonnegative slack of $x_j$ gives $\pi_j=d_jw_j\le\lambda_0$. If $x_j>0$, the slack of
$x_j$ is zero, so $\pi_j=\lambda_0+t_j\ge\lambda_0$. By Lemma~\ref{lem:cutoff}, condition~(b)
holds. This is valid also for scores equal to zero, including $d_j=0$.

(b) Let $(x,f)$ be a BNE in marginal coordinates. Since $0<r<n$, some target has $x_j>0$ and some has
$x_k<1$. By Lemma~\ref{lem:cutoff} we may choose
\[
  \lambda_0\in\Bigl[\max_{k:\,x_k<1}\pi_k,\ \min_{j:\,x_j>0}\pi_j\Bigr],
\]
and $\lambda_0\ge0$ because the scores are nonnegative and the first set is nonempty. Set
$\lambda=1+\lambda_0$, $u_i=\max_k(A_{ik}-b_{ik}x_k)$ and $t_j=(\pi_j-\lambda_0)_+$. We check the five
rows of~\eqref{app-dl:eq:lcp}. The slack of $f_{ij}$ is $u_i-(A_{ij}-b_{ij}x_j)\ge0$, and it vanishes
when $f_{ij}>0$ because $j$ is then a best target. The slack of $t_j$ is $1-x_j\ge0$; if $t_j>0$ then
$\pi_j>\lambda_0$, so $x_j=1$. The budget slack is $0$. The slack of $u_i$ is $0$. The slack of $x_j$ is
$\lambda_0+(\pi_j-\lambda_0)_+-\pi_j=(\lambda_0-\pi_j)_+\ge0$, and it vanishes when $x_j>0$ because
then $\pi_j\ge\lambda_0$. No tie-breaking between equally good targets is needed in either direction.
\end{proof}

\subsection{The matrix is \texorpdfstring{$P_0$}{P0}}\label{app-dl:sec:p0}

Split the variables into the first block $(f,t,\lambda)$ and the second block $(u,x)$. The slacks of the
first block depend only on $(u,x)$, and those of the second block only on $(f,t,\lambda)$. So
\[
  M=\begin{pmatrix}0&B^{\mathsf T}\\-D&0\end{pmatrix},
\]
where $B$ and $D$ are matrices with rows indexed by the second block (one \emph{type row} $i$ per type and
one \emph{target row} $j$ per target) and columns indexed by the first block. We name them after the
payoff data they carry: $B$ carries the slopes $b_{ij}$ and $D$ the defender gains $d_j$.
Reading off~\eqref{app-dl:eq:lcp}, their columns are as follows; all unspecified entries are zero.
\begin{itemize}[nosep]
  \item \emph{Attack column $ij$.} In $B$: entry $1$ in type row $i$ and $b_{ij}$ in target row $j$. In
        $D$: entry $1$ in type row $i$ and $d_j$ in target row $j$.
  \item \emph{Cap column $t_j$.} In both $B$ and $D$: a single entry $-1$ in target row $j$.
  \item \emph{Budget column $\lambda$.} In both $B$ and $D$: entry $-1$ in every target row, $0$ in every
        type row.
\end{itemize}
(For example, the slack of $x_j$ is $\lambda-1+t_j-\sum_id_jf_{ij}$; its row of $-D$ has $+1$ at
$\lambda$ and $t_j$ and $-d_j$ at each $f_{ij}$.) Attack and cap columns are \emph{sparse}; the budget
column is \emph{dense}.

\paragraph{The block determinant.} A principal submatrix of $M$ selects a set $F$ of first-block indices
and a set $R$ of second-block indices. It is
\[
  M_{F\cup R,\,F\cup R}=\begin{pmatrix}0&(B_{R,F})^{\mathsf T}\\-D_{R,F}&0\end{pmatrix},
\]
where $B_{R,F}$ and $D_{R,F}$ have rows $R$ and columns $F$. If $|F|\ne|R|$, the rank of this
matrix is at most $2\min(|F|,|R|)$, which is less than $|F|+|R|$, so the minor is $0$. If $|F|=|R|=k$, then
\begin{equation}\label{app-dl:eq:blockdet}
\begin{aligned}
  \det\begin{pmatrix}0&(B_{R,F})^{\mathsf T}\\-D_{R,F}&0\end{pmatrix}
  &=(-1)^{k^2}\det(B_{R,F})^{\mathsf T}\det(-D_{R,F})\\
  &=(-1)^{k^2+k}\det B_{R,F}\det D_{R,F}=\det B_{R,F}\det D_{R,F}.
\end{aligned}
\end{equation}
The factor $(-1)^{k^2}$ comes from swapping the two block-column groups, and $(-1)^k$ from removing the
sign of the lower block; together they cancel. So it suffices to prove
$\det B_{R,F}\det D_{R,F}\ge0$ for all equal-size $R,F$.

\begin{lemma}[$P_0$]\label{app-dl:lem:p0}
\statuslean
For all data with $b_{ij}\ge0$ and $d_j\ge0$, every principal minor of $M$ is nonnegative.
\end{lemma}

\begin{proof}
\emph{Step 0: strictly positive data.} First assume $b_{ij}>0$ and $d_j>0$ for all $i,j$. Fix $R,F$ with
$|R|=|F|$ and write $B'$, $D'$ for the following rescaled versions of $B_{R,F}$, $D_{R,F}$. Negate every
type row in \emph{both} matrices. This multiplies each of the two determinants by $(-1)^{\#\text{type
rows in }R}$, the same factor, so it does not change their product. Divide every target row $j$ of $D$
by $d_j>0$; this multiplies $\det D_{R,F}$ by a positive number. After these operations:
\begin{itemize}[nosep]
  \item an attack column $ij$ has entries $(-1,+1)$ in $D'$ and $(-1,b_{ij})$ in $B'$, at type row $i$
        and target row $j$ (restricted to those rows that lie in $R$);
  \item a cap column has a single entry, $-1$ in $B'$ and $-1/d_j<0$ in $D'$, at its target row $j$;
  \item the budget column has, at every selected target row $j$, the entry $-1$ in $B'$ and $-1/d_j<0$ in
        $D'$, and $0$ elsewhere. (Only the signs of these entries are used below.)
\end{itemize}
In particular $B'$ and $D'$ have the same zero pattern and the same sign at every nonzero entry. It
suffices to show $\det B'\cdot\det D'\ge0$.

A restricted sparse column has two entries of opposite sign (an attack column with both endpoints in
$R$), one nonzero entry (a cap column, or an attack column with one endpoint in $R$), or none. Let $G$ be
the graph on the vertex set $R$ whose edges are the selected two-entry columns; isolated vertices are
kept. $G$ is bipartite (type rows versus target rows) and simple. A one-entry column is attached to the
vertex of its entry. Every sparse column lies within one connected component of $G$.

\emph{Step 1: the cases in which the product is zero.} In each of the following cases
$\det B'\cdot\det D'=0$.
\begin{enumerate}[label=(Z\arabic*),nosep]
  \item Some selected sparse column has no entry in $R$. Then both matrices have a zero column.
  \item $G$ contains a cycle. Its edge columns in $D'$ are oriented incidence vectors $(-1,+1)$, and
        the alternating $\pm1$ combination around the cycle is zero. So $\det D'=0$.
  \item Some component $T$ of $G$ carries more selected sparse columns than it has vertices. These
        columns are supported on the $|T|$ rows of $T$, so they are linearly dependent in both matrices.
  \item The budget column is selected, and (after excluding (Z1)--(Z3)) the unique component $T_0$
        described in Step~3 below contains no target row. This case is treated in Step~3: both
        determinants vanish.
\end{enumerate}
From now on assume that (Z1)--(Z3) fail. Then $G$ is a forest, every component $T$ is a tree with
$|T|-1$ edge columns, and it carries at most one one-entry column.

\emph{Step 2: unique terms in trees.} Call a tree \emph{anchored} if it carries exactly one one-entry
column, and call the vertex of that entry its \emph{root}. In a square matrix whose columns are the
columns of anchored trees, the determinant expansion has exactly one nonzero term. Indeed, a nonzero
term is a bijection from columns to rows using nonzero entries only, and columns can only be matched
within their own tree. In a tree, a leaf other than the root has a nonzero entry only in its unique edge
column, so it must be matched to that edge. Removing the leaf and repeating, every edge is matched to its
endpoint farther from the root, and the one-entry column to the root. The same argument shows: if a
tree $T$ carries no one-entry column and one of its rows $\rho$ is deleted, the remaining $|T|-1$ rows of
$T$ and its $|T|-1$ edges have exactly one nonzero matching, with $\rho$ serving as root. Because $B'$
and $D'$ have the same zero pattern and the same entry signs, the unique term is given by the same
permutation in both matrices, with the same sign.

\emph{Case A: the budget column is not selected.} There are $|R|$ sparse columns. Each tree $T$ carries
$|T|-1$ or $|T|$ of them, and the total is $|R|=\sum_T|T|$. Hence every tree carries exactly $|T|$
columns, that is, every tree is anchored. By Step~2, $\det B'$ and $\det D'$ each consist of one nonzero
term, with the same sign. So $\det B'\cdot\det D'>0$.

\emph{Step 3 / Case B: the budget column is selected.} Put it last in both matrices. There are $|R|-1$
sparse columns, so by the same count exactly one tree $T_0$ carries no one-entry column (it is
\emph{unanchored}) and all others are anchored. For $X\in\{B',D'\}$ expand $\det X$ along the budget
column:
\[
  \det X=\sum_{v\in R}X_{v,\lambda}\,C^X_v=\sum_{j\in R\text{ target row}}X_{j,\lambda}\,C^X_j,
\]
where $C^X_v$ is the cofactor of position $(v,\lambda)$. The budget entry $X_{j,\lambda}$ of a target row is $-1$ in $B'$ and $-1/d_j$ in $D'$, negative in both; type rows have budget entry $0$. If $v$ lies in an anchored tree $T'$, deleting
row $v$ leaves the $|T'|$ columns of $T'$ supported on $|T'|-1$ rows, so $C^X_v=0$. Hence only target rows
of $T_0$ contribute.

\emph{Determinant-zero subcase (Z4).} If $T_0$ contains no target row (for example, $T_0$ is a single type
row), every term of the expansion vanishes, and $\det B'=\det D'=0$.

\emph{Nonzero subcase.} Now $T_0$ contains at least one target row. For $v\in T_0$, deleting row $v$ and
the budget column leaves anchored trees and the tree $T_0$ with root $v$, so by Step~2 the cofactor
$C^X_v$ is a single nonzero term. Moreover $C^{B'}_v$ and $C^{D'}_v$ have the same sign: they are given by
the same permutation, the entries have the same signs, and the position sign $(-1)^{v+\lambda}$ is the
same. Next, within one matrix all nonzero cofactors have one common sign. Indeed, the cofactor vector
$C^X=(C^X_v)_{v\in R}$ lies in the left kernel of every other column of $X$: for each sparse column $s$,
$\sum_vC^X_vX_{v,s}$ is the determinant of $X$ with the budget column replaced by $s$, which has a
repeated column. Applied to an edge column of $T_0$ with endpoints $i$ (type, entry $-1$) and $j$
(target, entry $\beta=1$ in $D'$ or $\beta=b_{ij}>0$ in $B'$), this gives $-C^X_i+\beta C^X_j=0$, so
$C^X_i$ and $C^X_j$ have the same sign. Since $T_0$ is connected and all its cofactors are nonzero, they
all have one common sign $\sigma_X$, and $\sigma_{B'}=\sigma_{D'}$ by the previous observation. We compare
signs only for these nonzero terms. Hence $\det X=\sum_{j\in T_0\text{ target}}X_{j,\lambda}C^X_j$ is a nonempty sum of terms of sign $-\sigma_X$, so $\det X$ is nonzero with sign $-\sigma_X$, the same for $X=B'$ and
$X=D'$. Thus $\det B'\cdot\det D'>0$.

Together with Step~1 and~\eqref{app-dl:eq:blockdet}, every principal minor of $M$ is nonnegative when
all $b_{ij},d_j$ are positive.

\emph{Step 4: zero coefficients, by continuity.} For arbitrary $b_{ij}\ge0$ and $d_j\ge0$, take a sequence
of strictly positive coefficient vectors converging to the given ones. Each principal minor of $M$ is a
polynomial in the coefficients, nonnegative along the sequence by Steps~0--3, hence nonnegative in the
limit. The input itself is not changed.
\end{proof}

\subsection{Perturbation, bound, gap and the main theorem}\label{app-dl:sec:main}

\begin{lemma}[Diagonal perturbation]\label{app-dl:lem:pmatrix}
\statuslean
For every rational $\delta>0$, $M+\delta I$ is a P-matrix. In particular
$M+\delta I$ has no sign-reversing vector, and $\mathrm{LCP}(q,M+\delta I)$ has at most one solution for every $q$.
\end{lemma}

\begin{proof}
For a nonempty principal index set $S$,
$\det(M+\delta I)_{SS}=\sum_{U\subseteq S}\delta^{|S|-|U|}\det M_{UU}\ge\delta^{|S|}>0$, because every
term is nonnegative by Lemma~\ref{app-dl:lem:p0} and the term $U=\emptyset$ equals $\delta^{|S|}$. The
rest is Proposition~\ref{prop:classical}(i) and the classical fact that a P-matrix has no nonzero $y$
with $y_k(My)_k\le0$ for all $k$~\cite{cps92} (see the remark after Proposition~\ref{prop:classical}).
\end{proof}

By Proposition~\ref{prop:classical}(i) the solution exists; this classical fact is not part of the Lean proof.

\paragraph{Constants.} Define
\[
  U=\max_{i,j}A_{ij},\quad W=1+mU,\quad D_*=\max_jd_j,\quad S=1+D_*W,\quad
  K=\bigl\lceil\max(1,U,W,1+S)\bigr\rceil .
\]
Note $U\ge1$ because $A_{ij}\ge1+b_{ij}$.

\begin{lemma}[Uniform bound]\label{app-dl:lem:bound}
\statuslean
For every $0<\delta\le1$, the solution $z=(f,t,\lambda,u,x)$ of $\mathrm{LCP}(q,M+\delta I)$ lies in
$[0,K]^N$.
\end{lemma}

\begin{proof}
The perturbed slack of each variable is its slack in~\eqref{app-dl:eq:lcp} plus $\delta$ times the
variable. \emph{Types.} If $u_i>0$, its slack vanishes: $\sum_jf_{ij}=p_i+\delta u_i>0$. On a positive
$f_{ij}$ the slack of $f_{ij}$ vanishes, so $u_i=A_{ij}-b_{ij}x_j-\delta f_{ij}\le U$. Thus $u_i\le U$ in
all cases. \emph{Flows.} The nonnegative type slacks give $\sum_jf_{ij}\le p_i+\delta u_i$, so
$\sum_{i,j}f_{ij}\le1+\delta mU\le W$; in particular every $f_{ij}$ and every $w_j$ is at most $W$.
\emph{Caps.} If $t_j>0$, then its slack vanishes, $x_j=1+\delta t_j>0$, so the slack of $x_j$ vanishes:
$t_j=1+d_jw_j-\lambda-\delta x_j\le1+D_*W=S$. \emph{Coverage.} The nonnegative cap slack gives
$x_j\le1+\delta t_j\le1+S$. \emph{Cutoff.} If $\lambda>0$, then $\sum_jx_j=r+\delta\lambda>0$, so some
$x_j>0$, and its vanishing slack gives $\lambda=1+d_jw_j-t_j-\delta x_j\le S$. All coordinates are at most
$K$.
\end{proof}

\paragraph{The pattern LP.} For a set $P\subseteq[N]$ (a \emph{pattern}), the \emph{pattern LP} in the
variables $z\in\mathbb R^N$ is
\[
  0\le z\le K,\qquad z_k=0\ (k\notin P),\qquad (q+Mz)_k\ge0\ (\text{all }k),\qquad (q+Mz)_k=0\ (k\in P).
\]
Every feasible point solves $\mathrm{LCP}(q,M)$: off $P$ the variable is zero, on $P$ the slack is zero.
This holds even if some coordinate in $P$ is zero at the feasible point.

\paragraph{The relaxed pattern LP.} In the variables $(z,\tau)\in\mathbb R^{N+1}$, minimize $\tau$ subject
to the following rows.
\begin{enumerate}[label=(R\arabic*),nosep]
  \item For every $k\in[N]$: $(q+Mz)_k+\tau\ge0$ (the relaxed slack inequality; coefficient of $\tau$
        is $+1$).
  \item For every $k\in P$: $-(q+Mz)_k+\tau\ge0$. Together with (R1) this replaces the equality
        $(q+Mz)_k=0$ by the two inequalities $-\tau\le(q+Mz)_k\le\tau$ (coefficient of $\tau$ is $+1$).
  \item Boxes, exact: $z_k\ge0$ and $z_k\le K$ for every $k$ (coefficient of $\tau$ is $0$).
  \item Pins, exact: $z_k\le0$ for every $k\notin P$, which with (R3) gives $z_k=0$ (coefficient of
        $\tau$ is $0$).
  \item $\tau\ge0$ and $\tau\le1$.
\end{enumerate}
Every row that can occur for any pattern belongs to the fixed list consisting of the $2N$ rows
$\pm(q+Mz)_k+\tau\ge0$, the $3N$ rows $z_k\ge0$, $z_k\le K$, $z_k\le0$, and the two rows of (R5). It is
written with the coefficient of $\tau$ inserted \emph{before} clearing denominators. Clear the
denominators of each of these $5N+2$ rows separately (multiply the row by the least common multiple of
the denominators of its coefficients and its right side). Let $C\ge1$ be an integer bounding the absolute
values of all resulting integer coefficients and right sides, and put
\[
  Q=\bigl((N+1)C\bigr)^{N+1} .
\]

\begin{lemma}[Gap]\label{app-dl:lem:gap}
\statuslean
For every pattern $P$ for which the relaxed pattern LP is feasible, its optimal value is either $0$ or at
least $1/Q$.
\end{lemma}

\begin{proof}
The feasible region is a nonempty polytope, bounded by (R3) and (R5), so the minimum is attained at a
vertex. A vertex of a polytope in $\mathbb R^{N+1}$ is the unique solution of at most $N+1$ linearly
independent tight rows; choose exactly $N+1$ of them, giving a nonsingular integer system $Gy=h$ with
entries bounded by $C$. By Cramer's rule $\tau=\det G_\tau/\det G$, where $G_\tau$ is $G$ with the
$\tau$-column replaced by $h$. If $\tau>0$, then $|\det G_\tau|\ge1$, so
$\tau\ge1/|\det G|\ge1/\bigl((N+1)!\,C^{N+1}\bigr)\ge1/Q$.
\end{proof}

\begin{theorem}[Exact BNE $\le_m$ P-LCP]\label{app-dl:thm:main}\label{thm:direct-lcp}
\statusrefereed
Put $\delta=1/(2KQ)$. The map $\Gamma\mapsto(q,M+\delta I)$ is computable in time polynomial in $L$, and
$M+\delta I$ is a P-matrix. If $0<r<n$ the instance has dimension $N=mn+m+2n+1$; if $r$ is $0$ or $n$,
the instance map returns the fixed one-dimensional instance described in
Section~\ref{app-dl:sec:lcp}. Let $z^*$ be the unique solution of $\mathrm{LCP}(q,M+\delta I)$ and let
$P=\{k:z^*_k>0\}$. Then the pattern LP of $P$ is feasible, each of its solutions solves
$\mathrm{LCP}(q,M)$, and hence (Lemma~\ref{app-dl:lem:formulation}) gives a BNE in marginal
coordinates. Together with the decoder of Lemma~\ref{app-dl:lem:decode}, this is a polynomial-time
many-one reduction from exact BNE search in Bayesian simple security games to P-LCP, using one oracle
query that satisfies the P-matrix promise.
\end{theorem}

\begin{proof}
\emph{The query is a P-matrix} by Lemma~\ref{app-dl:lem:pmatrix}.

\emph{Exact recovery.} Since $\delta\le1$, Lemma~\ref{app-dl:lem:bound} gives $z^*\in[0,K]^N$. The
perturbed slacks satisfy $(q+Mz^*)_k+\delta z^*_k\ge0$, so $(q+Mz^*)_k\ge-\delta K$ for every $k$. For
$k\in P$ the perturbed slack is zero, so $(q+Mz^*)_k=-\delta z^*_k\in[-\delta K,0]$. Thus
$(z^*,\tau=\delta K)$ satisfies (R1)--(R5), with $\tau=\delta K=1/(2Q)\le1$. The optimal value of the
relaxed pattern LP is at most $1/(2Q)<1/Q$, hence $0$ by Lemma~\ref{app-dl:lem:gap}. An optimal point
with $\tau=0$ is a feasible point of the pattern LP. Any solution of the pattern LP, computed by a
polynomial-time LP algorithm from the original payoffs, priors, budget and unit caps, solves
$\mathrm{LCP}(q,M)$. Coordinates in $P$ that become zero are harmless.

\emph{Size.} The matrix and $q$ have entries built from the input by a polynomial number of additions,
comparisons and maxima. $K$ is an integer with $\log K$ polynomial in $L$. Every coefficient and right
side in the list (R1)--(R5) is $0$, $\pm1$, $K$, an input number ($b_{ij}$, $d_j$, $p_i$, $r$) up to
sign, or a shifted intercept $A_{ij}$ (a sum of $1$ and three input numbers). So the least common
multiple of the denominators of a row, and the cleared integers, have polynomial bit length, and
$\log C$ is polynomial in $L$. Hence
\[
  \frac1\delta=2KQ,\qquad \log_2\frac1\delta=1+\log_2K+(N+1)\log_2\bigl((N+1)C\bigr),
\]
so $\delta$ is the reciprocal of an integer of bit length at most $2+\log_2K+(N+1)\log_2((N+1)C)$,
polynomial in $L$. The unique solution $z^*$ is rational of polynomial encoding length, since it solves
the nonsingular principal linear system of $M+\delta I$ on its support. The solution map (support, LP,
decoder) runs in polynomial time. Enumeration of minors and patterns is used only in the proofs, never by
the algorithms. Edge cases and zero-prior types are handled as in Section~\ref{app-dl:sec:lcp} and
Lemma~\ref{app-dl:lem:decode}.
\end{proof}

On the small games of Appendix~\ref{app:sources}, $1/\delta$ had between $56$ and $958$ bits.

\subsection{Returning the strategies}\label{app-dl:sec:decode}

\begin{lemma}[Decoder]\label{app-dl:lem:decode}
\statusrefereed
Let $(x,f)$ be a BNE in marginal coordinates of the positive-prior game. The following is computable in
polynomial time and is a BNE of the original explicit game.
\begin{enumerate}[label=(\alph*),nosep]
  \item Type $i$ with $p_i>0$ attacks $j$ with probability $f_{ij}/p_i$.
  \item Each zero-prior type plays one target maximizing $a_{ij}-b_{ij}x_j$.
  \item The defender plays the following mixed strategy on at most $n$ pure strategies, whose coverage
        marginals are exactly $x$.
\end{enumerate}
\end{lemma}

\begin{proof}
(a) is a probability distribution since $\sum_jf_{ij}=p_i$, and it uses only best targets. (b) A
zero-prior type has zero weight in the defender's expected payoff, so restoring it with a best response
does not change the scores $\pi_j$ or the defender's incentives.

(c) This is the construction of Lemma~\ref{lem:realize}, in detail. Let $c_0=0$ and $c_j=x_1+\dots+x_j$,
so $c_n=r$. Partition $[0,r)$ into the consecutive half-open intervals $I_j=[c_{j-1},c_j)$ of lengths
$x_j\le1$ (empty if $x_j=0$). For $\theta\in[0,1)$, the points $\theta,\theta+1,\dots,\theta+r-1$ all lie
in $[0,r)$, so each lies in exactly one interval. Let $S(\theta)$ be the set of targets whose intervals
contain one of these points. A half-open interval of length at most one contains at most one of the
points, which are spaced one apart. Hence the $r$ points hit $r$ distinct intervals, and $S(\theta)$ is an
$r$-subset of the targets, that is, a pure defender strategy. For $\theta$ uniform on $[0,1)$, target $j$
is selected exactly when $\theta$ lies in the image of $I_j$ modulo $1$; this image has measure
$|I_j|=x_j$ because $|I_j|\le1$. Endpoint choices form a set of measure zero and are resolved by the
half-open convention above. Let $0=\alpha_0<\alpha_1<\dots<\alpha_s<1$ be the distinct fractional parts of
$c_0,\dots,c_n$; since $c_0$ and $c_n$ have fractional part $0$, there are $s+1\le n$ of them. On each
piece $[\alpha_k,\alpha_{k+1})$ (with $\alpha_{s+1}=1$) no point $\theta+\ell$ crosses an interval
endpoint, so $S(\theta)$ is constant. Identify each piece's pure strategy as an $r$-subset, group pieces
with equal subsets, and give each subset the total length of its pieces. This is a rational mixed
strategy on at most $n$ pure strategies with marginals $x$. By Definition~\ref{app-dl:def:bne} and the
remark after it, the profile is a BNE.
\end{proof}

\begin{corollary}[Exact BNE is in UEOPL; ARRIVAL $\le_m$ P-LCP]\label{cor:direct-lcp}
\statusrefereed
Exact BNE search is in $\mathsf{UEOPL}$ (by Theorem~\ref{thm:direct-lcp} and \cite[Thm.~81]{fgms18}; the
query is a genuine P-matrix, so no violation exists). Composing Theorem~\ref{thm:A} with
Theorem~\ref{thm:direct-lcp} gives a second proof of ARRIVAL $\le_m$ P-LCP, through the game instead of through Theorem~\ref{plcp:thm:main}.
\end{corollary}

Since no polynomial-time algorithm for P-LCP is known,
Theorem~\ref{app-dl:thm:main} gives no polynomial-time algorithm for exact BNE, and since P-LCP is not
known to be $\mathsf{UEOPL}$-hard, it says nothing about hardness. The game may have many equilibria:
uniqueness is a property of the perturbed LCP, not of the game.

\section{One \texorpdfstring{$\ell_1$}{l1} contraction for the game}\label{sec:contraction}

\paragraph{Idea.} We want a map whose fixed point is (close to) an equilibrium and which contracts.
\begin{enumerate}[nosep]
\item \emph{Attackers.} Replace each type's best response by a soft version: project the payoff vector,
divided by a small temperature $\sigma$, onto the probability simplex. Positive flows then sit within
$\sigma$ of the best payoff, and the projection's Jacobian has zero column sums.
\item \emph{Defender.} Move coverage towards targets whose load exceeds the capacity
$\lambda/d_j$, by a small step $\eta$. One extra coordinate $q$ selects the cutoff $\lambda(q)$ and
tracks the budget.
\item \emph{Contraction.} The Jacobian of the combined update is nonnegative with column sums at most
one, so the map is $\ell_1$-nonexpansive; mixing in a constant point with weight $\delta$ makes it a
$(1-\delta)$-contraction.
\item \emph{Repair.} At the fixed point every equilibrium condition holds up to an error below a
threshold $g$. A Cramer-rule gap shows that then one pattern LP has an exact solution.
\end{enumerate}
Throughout this section $p_i>0$, $\sum_ip_i=1$ and $0<r<n$. Zero slopes and zero gains are allowed.

\paragraph{Intuition.} In the example of Figure~\ref{fig:example-market}, type~1 is exactly tied between
its two targets. The soft response of step~1 replaces the tie by a smooth split, which moves continuously
as the coverage moves; step~2 then pushes coverage towards the target whose load exceeds $\lambda/d_j$.
An exact fixed point of the combined map is only an \emph{approximate} equilibrium, off by about $\sigma$
in payoffs and $\kappa$ in the budget, and step~4 is what turns it into an exact one.

\subsection{Pattern LPs and the repair threshold}

A \emph{pattern} specifies a set of allowed type--target edges and, for every target, a label in
$\{\text{zero},\text{one},\text{free}\}$. The \emph{pattern LP} has variables $x,f,v,\lambda$ and
requires:
\begin{itemize}[nosep]
  \item the budget $\sum_jx_j=r$, the bounds $0\le x\le1$, and the row sums $\sum_jf_{ij}=p_i$;
  \item $f_{ij}=0$ off the allowed edges, and $v_i\ge U_{ij}(x)$ for all $(i,j)$, with equality on
        allowed edges;
  \item $x_j=0$ and $d_jw_j\le\lambda$ on zero targets, $x_j=1$ and $d_jw_j\ge\lambda$ on one targets,
        and $d_jw_j=\lambda$ on free targets.
\end{itemize}
It also has the \emph{outer boxes} $x,f\in[0,1]$, $0\le\lambda\le1+\max_jd_j$ and
$|v_i|\le2+\max_{i,j}(|a_{ij}|+b_{ij})$. All conditions are closed and linear.

\begin{lemma}[Soundness and uniform repair threshold]\label{lem:ueopl-pattern}
\statusrefereed
\begin{enumerate}[label=(\alph*),nosep]
  \item Every feasible point of a pattern LP is an exact BNE, even if some allowed flows vanish.
  \item Relax every pattern constraint, writing each equality as two inequalities, by a common
        additive slack $t\ge0$. Keep the outer boxes exact and minimize $t$. Write every row that can
        occur in any pattern LP with the coefficient of $t$ inserted \emph{before} clearing
        denominators. Then clear denominators row by row, let $H\ge1$ bound all resulting integer
        entries and right-hand sides, let $N=n+mn+m+2$ be the number of variables including $t$, and
        put $Q=(NH)^N$. For every pattern, the optimal slack is either $0$ or at least $1/Q$.
  \item Consequently, put $g=1/(4Q)$. If a point satisfies all constraints of a pattern to accuracy
        strictly better than $g$ and satisfies the outer boxes exactly, then the unrelaxed LP of
        this pattern is feasible, and solving it gives an exact BNE.
\end{enumerate}
$Q$ is computable in polynomial time without enumerating patterns, and $\log Q=N\log(NH)$ is
polynomial in $L$.
\end{lemma}

\begin{proof}
(a) A feasible point satisfies (i) and (ii) of Lemma~\ref{lem:cutoff}(b). Flows off the allowed
edges are zero, and allowed edges are best responses whether or not their flow is positive.
(b) The relaxed polyhedron contains no line, so it is pointed: every variable except $t$ is boxed, and
$t$ is bounded below by $0$, so no direction can be followed both ways. It is nonempty (keep the outer
boxes and take $t$ large), and $t\ge0$ bounds the objective, so a minimum is attained at a vertex. By Cramer's rule, $t^*$ at a vertex is an integer divided by a basis determinant
of absolute value at most $N!\,H^N\le Q$. (c) The point gives a feasible relaxed solution with
$t<g<1/Q$, so the optimum is $0$. Then (a) applies.
\end{proof}

The same argument applies if $t$ is in addition restricted to $[0,1]$: a positive minimum is still
attained at a vertex and is at least $1/Q$, and a witness with all residuals below $g$ forces the minimum
to be zero. This is the form used in Proposition~\ref{prop:sx-min}. The rest of the construction
produces, in polynomial time from one fixed point, a point that meets the hypothesis of
Lemma~\ref{lem:ueopl-pattern}(c). The budget need hold only approximately at that point.

\subsection{Constants}

Let $g=1/(4Q)$ be as in Lemma~\ref{lem:ueopl-pattern}. Put $b_{\max}=\max_{i,j}b_{ij}$,
$d'_j=d_j+g/16>0$ and $\bar d=\max_jd'_j$. Choose the rational constants
\begin{gather*}
  \sigma=\kappa=\frac g8,\qquad
  \eta=\min\Bigl(1,\ \frac{\sigma}{1+b_{\max}},\ \frac{\kappa}{4(1+\bar d)(1+\sum_j1/d'_j)}\Bigr),\\
  \delta=\min\Bigl(\frac18,\ \frac{\eta g}{8(1+\bar d)},\ \frac{g}{8(1+b_{\max})}\Bigr).
\end{gather*}
They satisfy
\begin{equation}\label{eq:ueopl-const}
  \eta\,\bar d\sum_j\frac1{d'_j}\le\frac\kappa4\le\kappa(1-2\delta),
\end{equation}
because $\delta\le1/8$. Each constant is a minimum or product of $O(1)$ quantities of bit length $O(n(\log Q+L))$; the factor $n$ comes from the sum $\sum_j1/d'_j$. This includes the case $d_j=0$, where $1/d'_j=16/g=64Q$. So all constants have
polynomial bit length. They are rationals but need not be dyadic; no step uses dyadic values, and one
may round $\eta$ and then $\delta$ down to powers of two if dyadic constants are wanted, since every
requirement below is an upper bound on $\eta$ or $\delta$. The constant $\delta$ may be exponentially
small in $L$. Let $a^0=(r/n)\mathbf 1\in[0,1]^n$, and write $C_{l,h}(u)=\min(h,\max(l,u))$ for clipping.

\subsection{The map}

\paragraph{Attackers: simplex projection.} At $x\in[0,1]^n$, let $y_i(x)$ be the Euclidean projection
of the vector $U_i(x)/\sigma=(U_{ij}(x)/\sigma)_j$ onto the probability simplex on the $n$ targets. Put
$f_{ij}(x)=p_iy_{ij}(x)$ and $w_j(x)=\sum_ip_iy_{ij}(x)$.

\begin{lemma}[Projection circuit and its pieces]\label{lem:ueopl-proj}
Let $u\in\mathbb R^n$ have entries sorted as $u_{(1)}\ge\dots\ge u_{(n)}$. Its projection onto the
simplex is
\[
  y_j=\max(0,u_j-\theta),\qquad \theta=\max_{1\le k\le n}\frac{\sum_{h\le k}u_{(h)}-1}{k}.
\]
So $y_i(x)$ is computed by a circuit with a comparator sorting network (min/max gates), prefix sums,
division by the integer constants $k$, and $\max$. Moreover:
\begin{enumerate}[label=(\alph*),nosep]
  \item if $y_{ij}(x)>0$ then $U_{ij}(x)\ge\max_kU_{ik}(x)-\sigma$;
  \item on a region where the positive set of row $i$ is $S$ with $|S|=k$,
        \[
          \frac{\partial y_{ij}}{\partial x_h}=\frac{b_{ih}}{\sigma}\Bigl(\frac1k-[j=h]\Bigr)\qquad(j,h\in S),
        \]
        and all other entries of this Jacobian vanish;
  \item hence each column of $\partial w/\partial x$ sums to zero, its off-diagonal entries are
        nonnegative, its diagonal is at least $-b_{\max}/\sigma$, and $x+\eta w(x)$ has a nonnegative
        Jacobian with all column sums equal to one.
\end{enumerate}
\end{lemma}

\begin{proof}
The formula for $\theta$ is the standard solution of $\sum_j\max(0,u_j-\theta)=1$. (a) Every positive
coordinate has $u_j=\theta+y_j$ with $0<y_j\le1$. The maximal coordinate is at most $\theta+1$, since
$y\le 1$. So the two differ by at most one in units of $1/\sigma$, that is, by at most $\sigma$ in
payoff. (b) On $S$, $\theta=(\sum_{h\in S}u_h-1)/k$ and $y_j=u_j-\theta$, with
$\partial u_j/\partial x_h=-(b_{ih}/\sigma)[j=h]$. For $h\notin S$, $\theta$ and the positive $y_j$ do
not depend on $x_h$, and $y_h=0$. (c) Each column of (b) sums to $(b_{ih}/\sigma)(1-1)=0$, and the
off-diagonal entries are nonnegative. Summing with weights $p_i$, the diagonal is at least
$-\sum_ip_ib_{ih}/\sigma\ge-b_{\max}/\sigma$. Since $\eta\le\sigma/(1+b_{\max})$, the diagonal of
$I+\eta\,\partial w/\partial x$ is at least $1-b_{\max}/(1+b_{\max})>0$, and its column sums are one.
\end{proof}

A zero slope $b_{ih}=0$ only zeroes the corresponding column in (b).

\paragraph{Defender: one cutoff coordinate.} Add one coordinate $q\in[0,1]$ and put
$s=C_{0,\kappa}(q)$. Let $\lambda(s)$ equal $\bar d$ for $s\le\delta\kappa$, decrease affinely to $0$ at
$s=(1-\delta)\kappa$, and equal $0$ for larger $s$. That is,
$\lambda(s)=C_{0,\bar d}\bigl(\bar d\,((1-\delta)\kappa-s)/((1-2\delta)\kappa)\bigr)$. Define
\begin{gather*}
  T_j(x,s)=C_{0,1}\Bigl(x_j+\eta\bigl(w_j(x)-\lambda(s)/d'_j\bigr)\Bigr)\quad(j=1,\dots,n),\\
  R(x,s)=C_{0,\kappa}\Bigl(s+\sum_jx_j-\sum_jT_j(x,s)\Bigr),\\
  F(x,q)=(1-\delta)\bigl(T(x,s),R(x,s)\bigr)+\delta\,(a^0,\kappa/2).
\end{gather*}
Because $T\in[0,1]^n$, $R\in[0,\kappa]\subseteq[0,1]$ and $(a^0,\kappa/2)\in[0,1]^{n+1}$, $F$ maps
$[0,1]^{n+1}$ into itself. Every division is by a rational constant. So $F$ is a LinearFIXP circuit in
the sense of~\cite{fgms18} (Section~\ref{sec:pre-classes}). It has $O(mn^2+n)$ gates, counting a
comparator sorting network of $O(n^2)$ comparators per type, and polynomially many constant bits.
Clipping $q$ to $s$ makes $F$ a circuit on the ordinary unit cube. The narrow interval for $s$ lets the
fixed point select the defender's cutoff while keeping the budget accurate enough for
Lemma~\ref{lem:ueopl-pattern}.

\subsection{Contraction}

\begin{lemma}[$\ell_1$ contraction]\label{lem:ueopl-contraction}
\statusrefereed
For all $u,u'\in[0,1]^{n+1}$, $\|F(u)-F(u')\|_1\le(1-\delta)\|u-u'\|_1$. In particular $F$ has a unique
fixed point in $[0,1]^{n+1}$. Zero slopes $b_{ij}=0$ and zero gains $d_j=0$ are allowed.
\end{lemma}

\paragraph{Proof roadmap.} (1) On one affine piece, bound the Jacobian column by column: the $x$
columns by Lemma~\ref{lem:ueopl-proj}(c), the $s$ column by~\eqref{eq:ueopl-const}, and the $R$ row by
construction; clipping only lowers column sums. (2) Extend from pieces to all pairs along segments.
(3) Check that zero slopes and gains change nothing.

\begin{proof}
\emph{One affine piece.} Fix an affine piece and treat $s$ for the moment as an independent input.
\begin{itemize}[nosep]
  \item \emph{The $x$ columns of $\partial T$.} On rows where $T_j$ is unclipped, they are the rows of
        $I+\eta\,\partial w/\partial x$. On clipped rows they are zero. By
        Lemma~\ref{lem:ueopl-proj}(c), they are nonnegative with column sums at most one.
  \item \emph{The $s$ column of $\partial T$.} On an unclipped row $j$ its entry is $-\eta\lambda'(s)/d'_j$;
        on a clipped row it is zero. Here
        $\lambda'(s)\in\{0,-\bar d/((1-2\delta)\kappa)\}$. So it is nonnegative, and
        by~\eqref{eq:ueopl-const} its sum is at most
        $\eta\bar d\sum_j(1/d'_j)/(\kappa(1-2\delta))\le1$.
  \item \emph{The row of $R$.} Before clipping $R$, its entry in each column, $x_h$ or $s$, is one minus
        that column's sum in $\partial T$. This is nonnegative by the two previous items. Hence the
        Jacobian of $(T,R)$ with respect to $(x,s)$ is nonnegative with all column sums equal to one.
\end{itemize}
Clipping $R$ zeroes its row. The chain rule through $s=C_{0,\kappa}(q)$ multiplies the $s$ column by $0$
or $1$. Both steps keep the entries nonnegative and can only lower column sums. The $\ell_1$ operator norm
is the largest absolute column sum, so on every affine piece the Jacobian of $(T,R)$ as a function of
$(x,q)$ has $\ell_1$ operator norm at most one, and that of $F$ at most $1-\delta$.

\emph{All pairs.} $F$ is continuous and piecewise linear. Subdivide the segment $[u,u']$ at its finitely
many crossings of piece boundaries and apply the piecewise bound to each subsegment, using limits at the
boundaries. In more detail, since $F$ is continuous and piecewise affine, the cube has a finite partition into
relatively open polyhedral cells on each of which $F$ agrees with an affine map; on every cell choose
one affine realization of $F$, whose linear part is one of the Jacobians above. (Refine the partition so that the positive set $S_i$ of every row is constant on each cell. Gate ties may persist on a whole cell, and then not every tied realization has a linear part from the list above. On each cell use the realization in which $\theta$ is computed from that cell's positive set, with $k=|S_i|$; Lemma~\ref{lem:ueopl-proj}(b) gives its linear part.) The segment $[u,u']$ meets finitely many cells
in finitely many subintervals. On each subinterval the directional derivative of $F$ along $u'-u$ is
given by the chosen affine realization, so it has $\ell_1$ norm at most $(1-\delta)\|u'-u\|_1$.
Integrating along the segment and using continuity of $F$ at cell boundaries gives the bound for all
pairs. Uniqueness and existence
of the fixed point follow from Banach's theorem, since $F$ maps the compact cube into itself.

\emph{Degenerate data.} A zero slope only zeroes Jacobian columns in Lemma~\ref{lem:ueopl-proj}(b). A
zero gain enters only through $d'_j=g/16>0$, and $1/d'_j=64Q$ keeps the constants polynomial.
\end{proof}

\subsection{Decoding the fixed point}

\begin{lemma}[Budget at the fixed point]\label{lem:ueopl-budget}
Let $(x,q)$ be the fixed point of $F$. Then $q\in[\delta\kappa/2,(1-\delta/2)\kappa]$, $s=q$, $R$ is
unclipped, and
\[
  \sum_jx_j+q=r+\frac\kappa2,\qquad \Bigl|\sum_jT_j-r\Bigr|\le\frac\kappa{2(1-\delta)}.
\]
\end{lemma}

\begin{proof}
The last coordinate gives $q=(1-\delta)R+\delta\kappa/2$ with $R\in[0,\kappa]$. This gives the
interval for $q$, and $s=q$ because $0\le q\le\kappa$. Suppose $R=0$. Then $q=\delta\kappa/2\le\delta\kappa$,
so $\lambda=\bar d$. Since $w_j\le1$ and $d'_j\le\bar d$, we get $w_j-\lambda/d'_j\le0$ and $T\le x$. The
argument of the clip defining $R$ is then at least $q>0$, a contradiction. Suppose $R=\kappa$. Then
$q=(1-\delta/2)\kappa\ge(1-\delta)\kappa$, so $\lambda=0$ and $T\ge x$. The argument is then at most
$q<\kappa$, again a contradiction. So $R=q+\sum_jx_j-\sum_jT_j$.

The $x$ coordinates give $x=(1-\delta)T+\delta a^0$, so $\sum_jx_j=(1-\delta)\sum_jT_j+\delta r$.
Substituting into $q=(1-\delta)R+\delta\kappa/2$ and cancelling gives $\sum_jx_j+q=r+\kappa/2$. Then
$\sum_jT_j-r=(\kappa/2-q)/(1-\delta)$, and the bound follows from $q\in[0,\kappa]$.
\end{proof}

The interval for $q$ gives $|\kappa/2-q|\le(1-\delta)\kappa/2$, which yields the sharper bound
$|\sum_jT_j-r|\le\kappa/2$. Either bound suffices below.

\begin{lemma}[Decoding]\label{lem:ueopl-decoding}
\statusrefereed
Let $(x,q)$ be the fixed point and $z=T(x,q)$. Build a pattern as follows. The allowed edges are the
pairs with $f_{ij}(x)=p_iy_{ij}(x)>0$. Each target gets the label zero, one or free according to
whether $z_j=0$, $z_j=1$ or $0<z_j<1$. Then the point $(z,\,f(x),\,v,\,\lambda(q))$ with
$v_i=\max_jU_{ij}(z)$ satisfies the outer boxes exactly and all constraints of this pattern with
residual strictly below $g$. Hence, by Lemma~\ref{lem:ueopl-pattern}(c), one LP with the original
coefficients returns an exact BNE.
\end{lemma}

\paragraph{Proof roadmap.} Bound the residual of each kind of pattern constraint at
$(z,f(x),v,\lambda(q))$: exact constraints have residual $0$; best responses are off by at most
$\sigma+2b_{\max}\delta$ (Lemma~\ref{lem:ueopl-proj}(a)); cutoff conditions by at most
$\bar d\delta/\eta+g/16$; the budget by Lemma~\ref{lem:ueopl-budget}. All are below $g$, and
Lemma~\ref{lem:ueopl-pattern}(c) finishes.

\begin{proof}
Since $x=(1-\delta)z+\delta a^0$ and both lie in $[0,1]^n$, $|x_j-z_j|\le\delta$.

\emph{Exact constraints.} The bounds and boxes hold. We have $z\in[0,1]^n$, $f\in[0,1]$,
$0\le\lambda\le\bar d\le1+\max_jd_j$, and $|v_i|\le\max_{i,j}(|a_{ij}|+b_{ij})$. The row sums are exact,
flows off the pattern are zero, $v_i\ge U_{ij}(z)$ holds for all pairs, and $z_j\in\{0,1\}$ holds exactly
on zero and one targets.

\emph{Best responses.} On an allowed edge, Lemma~\ref{lem:ueopl-proj}(a) gives a payoff at $x$ within
$\sigma$ of the row maximum. Moving from $x$ to $z$ changes every payoff by at most $b_{\max}\delta$. So
$0\le v_i-U_{ij}(z)\le\sigma+2b_{\max}\delta<g/8+g/4$, using $\delta\le g/(8(1+b_{\max}))$.

\emph{Cutoff.} If $z_j>0$, the argument of the clip defining $T_j$ is at least $z_j$. So
$\eta(w_j-\lambda/d'_j)\ge z_j-x_j\ge-\delta$, that is, $d'_jw_j\ge\lambda-\bar d\delta/\eta$. If
$z_j<1$, the reverse argument gives $d'_jw_j\le\lambda+\bar d\delta/\eta$. Replacing $d'_j$ by $d_j$
changes the score by at most $(g/16)w_j\le g/16$. So every cutoff residual is at most
$\bar d\delta/\eta+g/16\le g/8+g/16$, using $\delta\le\eta g/(8(1+\bar d))$.

\emph{Budget.} By Lemma~\ref{lem:ueopl-budget}, $|\sum_jz_j-r|\le\kappa/(2(1-\delta))\le(8/7)(g/16)<g$.

All residuals are strictly below $g$, and Lemma~\ref{lem:ueopl-pattern}(c) applies.
\end{proof}

The decoder evaluates the circuit exactly at the fixed point and solves one rational LP. Both take
polynomial time. Zero-prior types are then given best responses, and Lemma~\ref{lem:realize} produces
explicit mixed strategies if they are wanted. In one statement:

\begin{theorem}[The game's $\ell_1$ contraction]\label{thm:sx-game}
\statusrefereed
$F:[0,1]^{n+1}\to[0,1]^{n+1}$ satisfies $\|F(u)-F(u')\|_1\le(1-\delta)\|u-u'\|_1$, and from its unique
fixed point an exact BNE is computed in polynomial time by one pattern LP. This holds with zero slopes
and zero gains.
\end{theorem}

\begin{proof}
Lemma~\ref{lem:ueopl-contraction} gives the contraction and the unique fixed point.
Lemmas~\ref{lem:ueopl-budget} and~\ref{lem:ueopl-decoding} decode it: at the fixed point $R$ is
unclipped, $\sum_jx_j+q=r+\kappa/2$, every $|z_j-x_j|\le\delta$, the used attacker payoffs at $z$ are
within $\sigma+2b_{\max}\delta$ of their row maxima, the cutoff residuals are at most
$\bar d\delta/\eta+g/16$, $\lambda\le\bar d\le1+\max_jd_j$, and all residuals are below $g$; so the
pattern LP recorded at $z$ (positive flow support, labels zero, one or free, and $v_i=\max_jU_{ij}(z)$)
is feasible by Lemma~\ref{lem:ueopl-pattern} and its solution is an exact BNE.
\end{proof}

\section{Exact BNE search is in UEOPL}\label{sec:ueopl}

\paragraph{Idea.} PL-Contraction asks for the fixed point of a piecewise-linear contraction \emph{or} a
certificate that the map is not one. Our $F$ is genuinely a contraction, maps the cube into itself and
has a fixed point, so no certificate exists and the only solution is the fixed point.

\begin{theorem}[Exact BNE search is in UEOPL]\label{thm:ueopl-main}
\statusrefereed
Exact BNE search for Bayesian simple security games with rational data written in binary reduces in
polynomial time to computing the unique fixed point of the LinearFIXP circuit map
$F:[0,1]^{n+1}\to[0,1]^{n+1}$ constructed above. Zero slopes and zero gains are allowed. $F$ is a
$(1-\delta)$-contraction in $\ell_1$ for the explicit rational $\delta>0$ of polynomial bit length.
Read as the \textnormal{\textsc{PL-Contraction}} instance $(F,\ c=1-\delta,\ p=1)$, it has no violation
solutions, so its only solution is the fixed point. Hence exact BNE search is in $\mathsf{UEOPL}$.
\end{theorem}

\begin{proof}
Delete the zero-prior types. If $r\in\{0,n\}$, map the game to a fixed
trivial instance. The decoder returns the forced coverage $x=\mathbf 0$
or $x=\mathbf 1$ and chooses a maximizing target for each type.
At an interior budget restore each deleted type with a best response to the recovered coverage.
The instance map $\Gamma\mapsto(F,1-\delta,1)$ runs in polynomial time: $Q$, the constants and the
circuit all have polynomial size. Lemma~\ref{lem:ueopl-contraction} excludes CMV1, and CMV2 is excluded
because $F$ maps the cube into itself. CMV3 is excluded because $F$ is contracting and has a fixed point
in the cube (\cite[Lemmas~30--31 and the remark after Def.~32]{fgms18}). So the only solution is the fixed point. Theorem~\ref{thm:fgms} applies, and
Lemma~\ref{lem:ueopl-decoding} decodes the fixed point with one LP. The game itself may have many
equilibria. Uniqueness is a property of $F$, not of the game.
\end{proof}

\begin{remark}[CLS]\label{rem:ueopl-cls}
$\mathsf{UEOPL}\subseteq\mathsf{EOPL}=\mathsf{CLS}$~\cite{fgms18,gjhmprt22,fghs23}, so
Theorem~\ref{thm:ueopl-main} also puts exact BNE search in $\mathsf{EOPL}$ and in
$\mathsf{CLS}=\mathsf{PPAD}\cap\mathsf{PLS}$~\cite{fghs23}.
\end{remark}

\begin{remark}[Comparison with ARRIVAL]\label{rem:ueopl-arrival}
Theorem~\ref{thm:ueopl-main} gives exact BNE search the same $\mathsf{TFNP}$ class bound as the search version of ARRIVAL,
$\mathsf{UEOPL}$, hence $\mathsf{EOPL}$ and $\mathsf{CLS}$. It does so through an \emph{exact}
piecewise-linear $\ell_1$ contraction in dimension $n+1$, whereas Haslebacher's reduction for ARRIVAL
goes to an \emph{approximate} fixed point~\cite{haslebacher25}. Hardness for $\mathsf{CLS}$ or
$\mathsf{EOPL}$ under many-one reductions would now give $\mathsf{EOPL}\subseteq\mathsf{UEOPL}$. This
inclusion fails relative to an oracle, that is, for decision-tree many-one reductions~\cite[Thm.~5]{goos22}.
That is a real barrier to black-box reductions of that kind, but it does not rule out an ordinary
polynomial-time hardness reduction.
\end{remark}

\begin{remark}[No polynomial algorithm follows]\label{rem:ueopl-nopoly}
For initial error $E$, the generic contraction estimate bounds the error
after $k$ iterations by $(1-\delta)^k E$. The iteration count obtained
from this estimate can be exponential in $L$. Within a known choice of affine pieces the fixed point solves a rational linear
system. The difficulty is identifying the pieces when the solution lies on or near changes in the
attackers' choices. Daskalakis, Farina and Zhang compute $\varepsilon$-approximate fixed points of
$\ell_p$-nonexpansive maps in time polynomial in $d$ and $\log(1/\varepsilon)$ for even integers
$p$~\cite{dfz26}. This excludes $p=1$.  In the other direction, a generic polynomial-time $\ell_1$ solver would also put ARRIVAL in
$\mathsf P$ by~\cite{haslebacher25}. A targeted algorithm would have to exploit the special structure
here, including the simplex projections in the attacker update and the
separate cutoff coordinate in the defender update.
\end{remark}

\begin{remark}[Fixed dimension]\label{rem:ueopl-fixeddim}
\cite[Theorem~40]{fgms18} and Theorem~\ref{thm:ueopl-main} imply a
polynomial-time exact algorithm for each fixed number of targets.
Here is a direct finite algorithm for this game. Write $n$ for the
number of targets, $m$ for the number of types, and $L$ for the binary
input length. In the affine hyperplane $\sum_jx_j=r$, list the
$2n$ equations $x_j=0,1$ and the $m\binom n2$ equations
$U_{ij}(x)=U_{ik}(x)$ for all types and unordered pairs of targets.
For each choice of $n-1$ equations, solve them together with the budget
equation. Keep a unique solution in $[0,1]^n$. At each retained
coverage find the best attacks of each type by exact comparison.
Solve the LP for attack flows and a cutoff on this menu: use full
prior rows, put zero on other attacks, and impose the cutoff
inequalities of Lemma~\ref{lem:cutoff} at the zero, interior and full
coordinates. A feasible candidate gives a full equilibrium by
Lemma~\ref{lem:realize}. Zero-prior types can choose any best attack.

For completeness, choose an actual equilibrium $(x^*,f^*)$ and fix its
flow and cutoff. Keep every coverage bound tight at $x^*$ as an
equality and require every attack used by $f^*$ to remain a best
response, using closed inequalities. Together with the budget and box
these constraints give a nonempty bounded polytope. At every point of
it the same flow and cutoff still work. An originally interior target
has score exactly the cutoff, which is admissible even if the target
reaches a bound. At a vertex the active equations contain the budget
and $n-1$ other independent equations, all from the list. Thus the
enumeration finds at least one feasible candidate. This includes ties
and identically zero slopes. Even for $n=1$ the empty choice works.
There are at most $\binom{2n+m\binom n2}{n-1}$ candidates and each
requires polynomial bit work. For fixed $n$ this is polynomial in $L$.
The candidate bound has an exponent depending on $n$.
\end{remark}
\noindent{\small\itshape AI self-review, November 26, 2026 (UTC).
Checked the vertex enumeration and preservation of the fixed attack
flow when an interior coordinate reaches a bound.}\par

\section{A subexponential exact algorithm}\label{sec:subexp}

\paragraph{Idea.} The fixed point of a $d$-dimensional $\ell_1$ contraction can be found by splitting
the coordinates into $\sqrt d$ blocks of size $\sqrt d$. For a guess of the first block, solve the rest
recursively; then cut away, in the first block, an $\ell_1$ halfspace that cannot contain the fixed point.
Ghorbani, Lopes and Mnich showed how to choose the guesses deterministically so that each cut removes at
least a $\frac1{2(k+1)}$ fraction of a grid of block points; combined with slicing (Lemma~\ref{lem:sx-slice}), this gives roughly $(\text{bits})^{O(\sqrt d)}$ work. Applied to the
$(n+1)$-dimensional map of Section~\ref{sec:contraction}, this solves the game.

\begin{theorem}[Subexponential exact BNE]\label{thm:sx-main}
\statusrefereed
An exact BNE of an explicitly given Bayesian simple security game with $n$ targets and binary input
length $L\ge2$ can be computed deterministically in $L^{O(\sqrt n)}$ bit operations. Zero attacker
slopes $b_{ij}$, zero defender gains $d_j$, rational payoffs and rational priors are allowed. In
particular the problem is solvable in $2^{O(\sqrt L\log L)}$ bit operations, and in polynomial time for
every fixed number of targets.
\end{theorem}

Theorem~\ref{thm:sx-main} follows from Theorem~\ref{thm:sx-pl} applied to the contraction of
Theorem~\ref{thm:sx-game}; the deduction is at the end of Section~\ref{sec:sx-recon}. The algorithm gives
nothing new for ARRIVAL: the reduction of Chapter~\ref{ch:arrival} produces games with $O(n^2)$
targets and types and $\mathrm{poly}(n)$ bits, on which even the bound of Proposition~\ref{prop:sx-min}
below is $2^{O(n\log n)}$, slower than the known ARRIVAL algorithms (Section~\ref{sec:pre-arrival}).

\paragraph{Relation to known results.} Theorem~\ref{thm:sx-pl} below combines known tools: the deterministic $\ell_1$ deep-point block search of Ghorbani, Lopes and Mnich~\cite{glm26}, developed there for ARRIVAL, and a slicing lemma (Lemma~\ref{lem:sx-slice}). It is not a new technique, and \cite{glm26} do not state it; this section adds the slicing step for strict contractions, the error control in the recursion and the exact reconstruction. They guess blocks of $k$ coordinates, recurse on the rest, and search each block by cutting $\ell_1$ halfspaces
at deterministic $\frac1{2(k+1)}$-deep points; their block solver (their Theorem~1.2) uses
$O(k^3\log(kN+2))$ evaluations and $(2k\log(N+2))^{O(k)}$ additional time, and their overall running time
has the same shape as the bound $(\bar Q+1)^a\cdot A^{O(b)}$ of Section~\ref{sec:sx-time}. Their
cube-root exponent comes from a trade-off specific to ARRIVAL (a dominating set and a base case exponential in a distance parameter);
for a general $d$-dimensional map the same balance gives $\sqrt d$. Their halfspace search works on maps
with unit increments, which are $\ell_1$-nonexpansive, and needs a ``leaky'' map to handle ties; the
strict contraction gap here avoids ties. Haslebacher, Lill, Schnider and Weber~\cite{hlsw25} give the same
cut search for $\ell_1$ contractions with $O(d^2(\log(1/\varepsilon)+\log(1/(1-\lambda))))$ queries
($\lambda$ the contraction factor); as noted by Chen, Li and Yannakakis~\cite{cly26}, choosing each query
point by brute force takes $\Omega((1/\varepsilon)^k)$ time in a $k$-dimensional block, and the
deep-point lemma of~\cite{glm26} removes exactly this bottleneck. Feodorov and Haslebacher~\cite{fh26}
prove the $\ell_\infty$ analogue, $(\log(1/\varepsilon))^{O(\sqrt d\log d)}$ time, by pairing a
query-efficient algorithm with a decomposition theorem; for $\ell_1$ the deep-point lemma supplies the
first ingredient and Lemma~\ref{lem:sx-slice} the decomposition. \cite{cly26} recurse on $\ell_1$
contractions by fixing coordinates and solving the rest to a boosted precision (for constant block size),
and remark that a full $\ell_1$ decomposition was unclear in their nonexpansive setting;
Lemma~\ref{lem:sx-slice} uses the known gap $\delta$, which is allowed here because $\log(1/\delta)$ is
part of the input length. Haslebacher~\cite{haslebacher25} reduces G-ARRIVAL to $\ell_1$ contraction and
gives the coordinate-at-a-time recursive framework. The exact-recovery step (a vertex denominator bound
plus continued fractions) is standard for exact PL contraction problems in the style of~\cite{fgms18}.

\begin{theorem}[Exact PL $\ell_1$ contraction in $S^{O(\sqrt d)}$; after~\cite{glm26}]\label{thm:sx-pl}
\statusrefereed
Let $F:[0,1]^d\to[0,1]^d$ be computed by a rational piecewise-linear circuit whose gates are addition,
subtraction, multiplication by a rational constant, minimum and maximum. Suppose a rational
$0<\delta\le1$ certifies that $\|F(u)-F(u')\|_1\le(1-\delta)\|u-u'\|_1$ for all $u,u'$. Let $S\ge2$ be the
combined binary encoding length of the circuit (including its input and output coordinates) and
$\delta$. Then the exact, unique fixed point of $F$ can be computed deterministically in $S^{O(\sqrt d)}$
bit operations.
\end{theorem}

\paragraph{Intuition: block size one.} For $d=2$ and blocks of size one, the outer search is a search on
the first coordinate. At a test value $c$, solve the one-dimensional tail approximately, compute
$h=G_X(c,y)$, and discard the side of $c$ away from $h$, where by Lemma~\ref{lem:sx-protect} the fixed
point cannot lie. Testing a median of the surviving grid points halves the candidates each time; this is
the variant used in the end-to-end checks (Appendix~\ref{app:status}). For larger blocks the median is
replaced by the deep points of Lemma~\ref{lem:sx-deep}.

Throughout the proof, $\alpha=1-\delta$ and all norms are $\ell_1$ norms. The proof has four parts: a
geometric primitive (halfspaces and deep points), slicing, one recursive call with a protected grid
point, and exact reconstruction.

\subsection{The geometric primitive}

For a center $c\in\mathbb R^k$ and a direction $v\ne0$, the \emph{$\ell_1$ halfspace} is
\[
  \mathcal H(c,v)=\{z:\ \|z-c\|_1\le\|z-(c-\epsilon v)\|_1\ \text{for every }\epsilon>0\}.
\]
Positive rescaling of $v$ does not change $\mathcal H(c,v)$.

\begin{lemma}[Sign formula]\label{lem:sx-sign}
Let $\sigma_i=\operatorname{sgn}(z_i-c_i)$. Then $z\in\mathcal H(c,v)$ if and only if
\[
  \sum_{i:\sigma_i\ne0}\sigma_iv_i+\sum_{i:\sigma_i=0}|v_i|\ \ge\ 0 .
\]
\end{lemma}

\begin{proof}
The left side is the right derivative at $\epsilon=0$ of the convex PL function
$\epsilon\mapsto\|z-c+\epsilon v\|_1$. A convex function is at least its value at $0$ for all $\epsilon>0$
if and only if its right derivative at $0$ is nonnegative.
\end{proof}

This tie convention (it matters at grid cuts and on singleton box sides) is the same as
in~\cite{glm26} (their Definition~3.2 and Lemma~7.1), which take the definition from~\cite{hlsw25}. A
point $w$ is \emph{$\beta'$-deep} in a finite set $C\subset\mathbb R^k$ if
$|C\cap\mathcal H(w,v)|\ge\beta'|C|$ for every $v\ne0$.

\begin{lemma}[Deep lattice point~\cite{glm26}]\label{lem:sx-deep}
\statuscited
Let $\beta\in\mathbb N$ and let $C\subseteq[0,2^\beta]^k$ be a nonempty disjoint union of $p$ integer
boxes, given by their corners; $|C|$ counts its integer points. One can compute deterministically an
integer point $w$ in the coordinate bounding box of $C$ that is $\frac1{2(k+1)}$-deep in $C$, in time
$(2k)^{O(k)}\cdot p\cdot\mathrm{poly}(\beta,\log(p+1))$.
\end{lemma}

This is the deep-point lemma of~\cite[\S8.4]{glm26}. Its guarantee applies to arbitrary finite unions
of such boxes, including lower-dimensional ones, and the bounding-box guarantee keeps every queried
center inside the domain even if it lies in a previously discarded part. We use the lemma independently
of the switch-graph algorithms of~\cite{glm26}; its centerpoint construction builds on the finite-point $\ell_1$ centerpoint theorem of~\cite{hlsw25} (Theorem~3.18 in its PDF, Theorem~3.20 in the arXiv HTML).

\subsection{Slicing preserves the contraction gap}

A contraction $G$ of a closed cube with fixed point $u^*$ satisfies the a posteriori bound
\begin{equation}\label{eq:sx-apost}
  \|u-u^*\|_1\le\|u-G(u)\|_1/\delta,
\end{equation}
since $\|u-u^*\|\le\|u-G(u)\|+\|G(u)-G(u^*)\|\le\|u-G(u)\|+\alpha\|u-u^*\|$.

A recursive call works with a \emph{restricted map} $G$: it freezes a prefix of the original input
coordinates, evaluates the original circuit, and keeps the remaining outputs. Fixing the same prefix in
both arguments and deleting outputs preserves the bound $\alpha$.

\begin{lemma}[Slicing]\label{lem:sx-slice}
\statusrefereed
Let $G=(G_X,G_Y)$ be an $\alpha$-contraction of a product of cubes with coordinates split into a block
$x$ of size $k\ge1$ and a tail $y$. For each $x$ the map $y\mapsto G_Y(x,y)$ is an $\alpha$-contraction;
let $y(x)$ be its fixed point and set $T(x)=G_X(x,y(x))$. Then $T$ is an $\alpha$-contraction, and if
$x^*$ is its fixed point, $(x^*,y(x^*))$ is the unique fixed point of $G$. For an empty tail, $T=G$.
\end{lemma}

\begin{proof}
By the fixed-point equations and additivity of the $\ell_1$ norm over coordinate blocks,
\[
  \|T(x)-T(x')\|_1+\|y(x)-y(x')\|_1=\|G(x,y(x))-G(x',y(x'))\|_1\le\alpha\bigl(\|x-x'\|_1+\|y(x)-y(x')\|_1\bigr).
\]
Rearranging, $\|T(x)-T(x')\|_1\le\alpha\|x-x'\|_1-(1-\alpha)\|y(x)-y(x')\|_1\le\alpha\|x-x'\|_1$. The
lifting statement is immediate from the definitions and uniqueness.
\end{proof}

\subsection{One call: precision and a protected grid point}

\paragraph{The call.} Given a restricted map $G$ and a rational tolerance $0<\tau\le1$, a call returns
a dyadic point of residual $\|u-G(u)\|_1\le\tau$. Split the free coordinates into a block $x$ of size $k$
and a tail. At a tested block point $c$, recursively solve the tail map $G_Y(c,\cdot)$ to residual
\[
  \theta=\delta^2\tau/16
\]
(with no tail, use the empty vector and residual zero). Call the result $y$ and evaluate $h=G_X(c,y)$
exactly with the original circuit. By~\eqref{eq:sx-apost} and contraction,
\[
  \|y-y(c)\|_1\le\delta\tau/16,\qquad \|h-T(c)\|_1\le e,\qquad e:=\delta\tau/16.
\]
If $\|h-c\|_1\le\tau/2$, return $(c,y)$; its full residual is at most $\tau/2+\theta\le\tau$. Otherwise
cut with the halfspace $\mathcal H(c,c-h)$; the direction $c-h$ is nonzero, rational, and computed
without rounding.

\paragraph{The grid.} Let $\beta=\lceil\log_2(16k/(\delta\tau))\rceil$, $N=2^\beta$, and search the grid
$\{0,1/N,\dots,1\}^k$. Initially all grid points are candidates; after a failed test at $c$, remove
$\mathcal H(c,c-h)$ from the candidates.

\begin{lemma}[Protected grid point]\label{lem:sx-protect}
Let $x^*$ be the fixed point of $T$ and let $z$ be a grid point obtained by rounding $x^*$
coordinatewise, so that $\rho:=\|z-x^*\|_1\le k/N\le e$. Then $z$ is never removed, even though the cuts
use approximate tail solutions.
\end{lemma}

\begin{proof}
Suppose $z\in\mathcal H(c,c-h)$ after a failed test at $c$. Taking $\epsilon=1$ in the definition gives
$\|z-c\|_1\le\|z-h\|_1$. Let $A=\|c-x^*\|_1$. Since $T(x^*)=x^*$,
\[
  A\le\rho+\|z-c\|_1\le\rho+\|z-h\|_1\le2\rho+\|h-x^*\|_1\le2\rho+\|h-T(c)\|_1+\|T(c)-T(x^*)\|_1\le2\rho+e+\alpha A .
\]
Hence $A\le(2\rho+e)/\delta\le3e/\delta=3\tau/16$, and
\[
  \|h-c\|_1\le\|h-T(c)\|_1+\|T(c)-x^*\|_1+\|x^*-c\|_1\le e+(1+\alpha)A\le\frac{\delta\tau}{16}+\frac{6\tau}{16}
  \le\frac{7\tau}{16}<\frac\tau2,
\]
contradicting the failed test.
\end{proof}

The algorithm never needs to identify $z$.

\paragraph{Implementing the cuts.} Store candidates in integer coordinates scaled by $N$. Each round
applies Lemma~\ref{lem:sx-deep} to the current candidate boxes, giving $w$, and tests $c=w/N$; positive
scaling preserves halfspace membership. A failed test removes at least a $\frac1{2(k+1)}$ fraction of the
candidates. The initial count is $(N+1)^k$ and $z$ always survives, so after $O(k^2\log(N+1))$ tests some
test succeeds. The argument allows $c$ outside the current candidate set and uses only that $c$ lies in
the cube.

After $r'$ cuts, divide each coordinate interval $[0,N]$ at every previous center coordinate, using a
singleton for the center and the nonempty gaps on either side. This gives at most $2r'+1$ intervals per
coordinate (also when centers coincide or lie on the boundary), hence at most $(2r'+1)^k$ disjoint
product boxes, and on each box every sign vector $\operatorname{sgn}(z-c)$ is constant for \emph{each}
previous center. To maintain the surviving union, test one corner of each product box against each cut
with Lemma~\ref{lem:sx-sign}. This is polynomial work per box, round, coordinate and bit of the
directions; no coordinate's $N+1$ values are ever enumerated. Frozen coordinates are absent from the
grid, and a coordinate collapsed to a singleton is handled by the zero-sign term of the formula.

\subsection{Recursion and running time}\label{sec:sx-time}

\begin{figure}[t]
\centering
\begin{tikzpicture}[
  nd/.style={draw,rounded corners,font=\scriptsize,align=center,inner sep=2pt},
  ar/.style={-{Stealth[length=1.6mm]}}]
  \node[nd,text width=5.6cm] (r) at (0,0) {level 1: block $x^{(1)}$ ($b$ coordinates), grid of spacing $2^{-\beta}$;
    test deep points $c_1,c_2,\dots$ (Lemma~\ref{lem:sx-deep}), cut $\mathcal H(c,c-h)$ after each failure};
  \foreach \i/\xx in {1/-4.2,2/-1.4,3/1.4} {
    \node[nd,text width=2.3cm] (c\i) at (\xx,-1.9) {tail solve at $c_\i$:\\level 2, tolerance $\theta=\delta^2\tau/16$};
    \draw[ar] (r) -- (c\i);
  }
  \node[font=\small] (dots) at (3.9,-1.9) {$\cdots$};
  \node[font=\scriptsize,align=left,anchor=west] at (4.4,-1.9) {at most $\bar Q=O(b^2(P+1))$\\tests per node};
  \foreach \i in {1,2,3} {
    \draw[ar] (c\i.south) -- ++(-0.5,-0.8);
    \draw[ar] (c\i.south) -- ++(0,-0.8);
    \draw[ar] (c\i.south) -- ++(0.5,-0.8);
  }
  \node[font=\scriptsize,align=center] at (-1.4,-3.6) {levels $3,\dots,a=\lceil d/b\rceil$: same pattern};
  \node[nd,text width=7.5cm] (leaf) at (-0.3,-4.4) {last level: block only (empty tail), exact evaluation of the circuit};
\end{tikzpicture}
\caption{The block recursion behind Theorem~\ref{thm:sx-pl}. A node searches one block of $b$ coordinates
by deep-point cuts; every test point $c$ triggers a recursive solve of the remaining coordinates, which by
Lemma~\ref{lem:sx-slice} is again an $\alpha$-contraction, and Lemma~\ref{lem:sx-protect} guarantees that
the inexact tail solutions never cut off the rounded fixed point. Accounting: each node makes at most
$\bar Q=O(b^2(P+1))$ tests, so the tree has $(\bar Q+1)^a$ nodes up to a constant; the deep-point
computations are local to a node and cost $A^{O(b)}W_*$ with $A=2+b+d+P$. The total is $A^{O(a+b)}W_*$;
with $b=\lceil\sqrt d\rceil$, so $a=O(\sqrt d)$, and $P$, $W_*$ polynomial in $S$, this is
$S^{O(\sqrt d)}$.}
\label{fig:sx-tree}
\end{figure}

The recursion is drawn in Figure~\ref{fig:sx-tree}.

Fix a block size $b$ with $1\le b\le d$ and let $a=\lceil d/b\rceil$ be the maximum number of nonempty
levels (the last block may be smaller). A call at depth $s$ has target $\tau_s=\tau(\delta^2/16)^s$. In
this subsection put
\[
  P=4+\lceil\log_2(16d/(\delta\tau))\rceil+a\lceil\log_2(16/\delta^2)\rceil .
\]
Every block exponent $\beta$ is at most $P$, and every dyadic coordinate produced anywhere in the
recursion has a denominator dividing $2^P$. Each $\tau_s$ is stored exactly with
$O(\mathrm{bits}(\tau)+a(\mathrm{bits}(\delta)+1))$ bits, and $\beta$ is found by integer comparisons with
powers of two.

\paragraph{Circuit evaluations.} For dyadic inputs with at most $P$ fractional bits, process the gates
in topological order. Every intermediate denominator divides $2^P$ times the product of the
denominators of the rational constant leaves and of the scalar-multiplication constants (each occurrence
counted once): addition and subtraction preserve a common denominator, a scalar gate adds the
denominator of its constant, and $\min$/$\max$ select an earlier value. Every magnitude is at most the
product of a factor $2$ per addition or subtraction, a factor $\max(1,|\text{constant}|)$ per scalar
gate, and the largest absolute leaf, so its logarithm is polynomial in $S$. Exact evaluations, residual
tests, cut directions and sign sums therefore use numbers of $\mathrm{poly}(S)+O(P)$ bits. A frozen prefix
never creates a new circuit.

\paragraph{Running time.} Each call makes at most $\bar Q=O(b^2(P+1))$ recursive tail calls, so after any
round the candidates form at most $(2\bar Q+1)^b$ boxes, and each application of
Lemma~\ref{lem:sx-deep} costs $(2b)^{O(b)}(2\bar Q+1)^b\,\mathrm{poly}(P,b\log(\bar Q+1))$; the cut
bookkeeping is of the same order up to polynomial factors in the direction bit lengths. With
$A=2+b+d+P$ and $W_*$ a fixed polynomial in $S$, $P$ and the tolerance encodings, the local work per node
is $A^{O(b)}W_*$. The recursion tree has $(\bar Q+1)^a$ nonempty nodes up to an absolute constant, so the
total work is
\[
  A^{O(a+b)}\,W_* .
\]
Only the polynomial number of tail evaluations branches at each level; the more expensive deep-point
computation is local to each node. Taking $b=\lceil\sqrt d\rceil$ gives $(d+P+2)^{O(\sqrt d)}W_*$ bit
operations for residual $\tau$. These bounds hold up to constants.

\subsection{Exact reconstruction}\label{sec:sx-recon}

\begin{lemma}[Denominator bound]\label{lem:sx-denom}
Introduce variables for the cube coordinates and for all gate values. For a choice of one active
argument at every $\min$ and $\max$ gate, impose the gate equalities, the weak inequalities that validate
the choices, the affine gate equalities, the cube bounds, and the equations identifying input
coordinates with output gates. Clear denominators row by row, let $H\ge1$ bound all resulting integer
entries and right sides, let $M$ be the number of variables, and put $D=(MH)^M$. Then every coordinate of
the fixed point $u^*$ of $F$ is a fraction with denominator at most $D$, and $\log D$ is polynomial in
$S$.
\end{lemma}

\paragraph{Proof roadmap.} Fix the branch pattern active at $u^*$. Its lifted polyhedron is a single point,
so that point is a vertex, and Cramer's rule bounds its denominators.

\begin{proof}
The menu of all possible rows has polynomial size, so $\log H$ and $M$ are polynomially bounded without
enumerating branch choices ($H$ itself can be exponential). Choose a branch pattern that occurs at $u^*$, resolving ties arbitrarily. Its
lifted polyhedron is a singleton: every feasible point projects to a fixed point of $F$, which is unique,
and topological evaluation then fixes every gate value. The active rows at this point span all $M$
coordinates; otherwise a nonzero common null vector could be followed a small distance in either
direction while keeping the strict slack of the inactive rows, contradicting the singleton property.
Cramer's rule on $M$ independent active rows bounds the denominator of every coordinate by $D$. This
covers branch regions of lower dimension and singular affine pieces.
\end{proof}

\begin{proof}[Proof of Theorem~\ref{thm:sx-pl}]
Run the recursion on $G=F$ with root tolerance $\tau=2^{-B}$, $B$ minimal with $\tau\le\delta/(8D^2)$;
then $B$, $P$ and all tolerance encodings are polynomial in $S$, so the running time of
Section~\ref{sec:sx-time} is $S^{O(\sqrt d)}$. By~\eqref{eq:sx-apost} the returned dyadic vector $u$
satisfies $\|u-u^*\|_1\le1/(8D^2)$. In each coordinate there is at most one fraction with denominator at
most $D$ within this distance, because two distinct such fractions are at least $1/D^2$ apart. The true
coordinate $p/q$, $q\le D$, satisfies $|u_i-p/q|<1/(2q^2)$, so by Legendre's criterion it is a convergent
of the continued fraction of $u_i$. List the convergents of $u_i$ (including $u_i$ itself) and pick the
one with denominator at most $D$ within distance $1/(8D^2)$. A final exact evaluation of $F$ verifies the
recovered point.
\end{proof}

\begin{proof}[Proof of Theorem~\ref{thm:sx-main}]
For $r\in\{0,n\}$ the answer is immediate. Otherwise apply Theorem~\ref{thm:sx-pl} to the map $F$ of
Theorem~\ref{thm:sx-game}: its dimension is $n+1$ and its encoding length $S$, including
$\log(1/\delta)$, is polynomial in $L$, so the fixed point is found in $L^{O(\sqrt{n+1})}=L^{O(\sqrt n)}$
bit operations. Decode it with one LP and convert to mixed strategies (Lemma~\ref{lem:realize}). Since
$n\le L$, this is $2^{O(\sqrt L\log L)}$.
\end{proof}

\subsection{A bound in the smaller dimension}

\paragraph{Idea.} If there are fewer types than targets, use the game-to-market reduction of
Chapter~\ref{ch:market}, transpose each market call so that types become goods, and turn the call back
into a game with about $2m$ targets. Theorem~\ref{thm:sx-main} then solves each call in
$L^{O(\sqrt m)}$.

\providecommand{\Xlo}{\underline{X}}\providecommand{\Xhi}{\overline{X}}
\paragraph{Setting and notation.} In this subsection a \emph{positive game} is a game with $p_i>0$ for
all $i$, $\sum_ip_i=1$, $b_{ij}>0$, $d_j>0$ and $0<r<n$. The letters $P$, $v$, $f$ have their meaning
in the market (M1) of Definition~\ref{def:m1} (prices, utilities, flows); the precision $P$ and the
level count $a$ of Section~\ref{sec:sx-time} do not occur here. For a rational number, its
\emph{binary length} is the total number of bits of its numerator and denominator.

\paragraph{Box queries.} Fix a positive game. A \emph{box query} is a triple $(c,l,h)$ of rational
vectors in $\mathbb Q^n$ with $c_j>0$ and $0\le l_j\le h_j\le1$ for all $j$. An \emph{answer} to it is an
equilibrium $(x,f)$ of the box market of Definition~\ref{def:mk-box} with the types as buyers
(masses $p_i$, payoffs $U_{ij}(x)=a_{ij}-b_{ij}x_j$), capacities $c$ and price box $[l,h]$; that is,
$l\le x\le h$, $f\ge0$, $\sum_jf_{ij}=p_i$ for all $i$, $f_{ij}>0$ only if $U_{ij}(x)=V_i(x)$, and, with
$w_j=\sum_if_{ij}$,
\begin{equation}\label{eq:sxm-box}
x_j>l_j\ \Longrightarrow\ w_j\ge c_j,\qquad x_j<h_j\ \Longrightarrow\ w_j\le c_j .
\end{equation}
A coordinate with $l_j=h_j$ is \emph{pinned}; \eqref{eq:sxm-box} is then vacuous for it. The game data
$(p,a,b)$ are part of every query.

\begin{lemma}[Interface of the game-to-market reduction]\label{lem:sxm-interface}
Let a positive game of binary length $L$ be given. The procedure of Lemmas~\ref{lem:mk-bisect}
and~\ref{lem:mk-narrow} (cutoff bisection, then narrowing with canonicalization and the final price LP)
has the following properties.
\begin{enumerate}[nosep,label=(\alph*)]
\item Every query is a box query $(c,l,h)$ with $c_j=\lambda/d_j$ for a rational $\lambda>0$. In the
bisection $l=\mathbf0$ and $h=\mathbf1$. In the narrowing $l=\Xlo$ and $h=\Xhi$ for the current endpoints
$\Xlo\le\Xhi$, except that $l_k=h_k=t$ at the pinned coordinate $k$, where $\Xlo_k\le t\le\Xhi_k$. If
$\lambda_0=0$, no query is made in the narrowing.
\item It makes at most $1+\lceil\log_2(DQ_1^2)\rceil+n\lceil\log_{3/2}(1/\eta)\rceil$ queries, a number
bounded by a fixed polynomial in $L$, and every query has binary length at most a fixed polynomial in
$L$.
\item It accepts \emph{any} answer to each query: whatever answers are returned, it stops with an exact
BNE $(x,f)$ of the game. Given the answers, it is deterministic.
\item If every answer has binary length at most a fixed polynomial in $L$, then all its work other than
answering the queries takes time polynomial in $L$.
\end{enumerate}
\end{lemma}

\begin{proof}
(a) The bisection queries box equilibria on $[0,1]^n$ at the midpoint cutoffs, which are positive
rationals (Lemma~\ref{lem:mk-bisect}). The narrowing queries the box market on $[\Xlo,\Xhi]$ with $k$
pinned to a grid point $t$ of the middle third of $[\Xlo_k,\Xhi_k]$, at the capacities
$c=\lambda_0/d$, and at $\lambda_0=0$ it uses the upper corner of the pinned sub-box instead
(Lemma~\ref{lem:mk-narrow}). The bracket satisfies $\Xlo\le\Xhi$ throughout, so $l\le h$, and
$0\le l\le h\le1$ because all endpoints lie in $[0,1]^n$. Since $d_j>0$ and $\lambda>0$, every $c_j>0$.

(b) The bisection makes at most $1+\lceil\log_2(DQ_1^2)\rceil$ queries (Lemma~\ref{lem:mk-bisect}). In
the narrowing, each query on coordinate $k$ shrinks the width $\Xhi_k-\Xlo_k$ by a factor at most $2/3$, no
width grows, every width starts at most $1$, and $k$ is queried only while its width exceeds $\eta$
(Lemma~\ref{lem:mk-narrow}). Hence $k$ is queried at most $\lceil\log_{3/2}(1/\eta)\rceil$ times. Since
$\log Q_1$, $\log Q_2$ and $\log(1/\eta)=\log(4nQ_2)$ are polynomial in $L$, so is the total. Bisection
queries have polynomial length by Lemma~\ref{lem:mk-bisect} (the midpoints and $d$), and narrowing
queries by Lemma~\ref{lem:mk-bits}.

(c) Lemmas~\ref{lem:mk-bisect} and~\ref{lem:mk-narrow} are stated for \emph{any} box equilibrium
returned by a query, and Lemma~\ref{lem:mk-pin} for any equilibrium of the pinned sub-box. A returned
point of total $r$ is a BNE by Lemma~\ref{lem:mk-box-cutoff}; otherwise the price LP of
Lemma~\ref{lem:mk-narrow} gives one. All other steps are rational linear programs and arithmetic, and
every choice that Section~\ref{sec:mk-g2m} leaves open is fixed by a rule: rounds process the
coordinates in index order and skip those of width at most $\eta$; $\lambda_0$ is the larger of the two
left ends; $t$ is the smallest grid point in the middle third; and every linear program is solved by one
fixed deterministic algorithm that returns a vertex.

(d) Proposition~\ref{prop:mk-g2m} proves this when the box answers come from (M1) answers through
Lemma~\ref{lem:mk-box-m1}; its proof uses only that the box answers have polynomial length. In detail,
an answer enters the procedure only through its total $\sum_jx_j$, its pattern (positive edges and
labels), and canonicalization, whose linear programs use only that pattern, the original data and the
current endpoints and pin; an answer of total $r$ also enters as the returned BNE. All linear programs of the procedure have rows of polynomial length by
Lemmas~\ref{lem:mk-bisect} and~\ref{lem:mk-bits}, and there are polynomially many of them by (b).
\end{proof}

\begin{lemma}[Answers of polynomial length]\label{lem:sxm-vertex}
Let $(c,l,h)$ be a box query of a positive game, let $\sigma$ bound the binary length of every number
among $p,a,b,c,l,h$, and put $N_{\mathrm v}=n+mn$. From any answer $(x,f)$ one computes, in time
polynomial in $n$, $m$, $\sigma$ and the binary length of $(x,f)$, an answer $(x',f')$ each of whose
entries has binary length at most $2N_{\mathrm v}(\log_2N_{\mathrm v}+4\sigma+2)$.
\end{lemma}

\begin{proof}
Let $A=\{(i,j):f_{ij}>0\}$, and give each good $j$ the label \emph{pinned} if $l_j=h_j$, and otherwise
\emph{floor}, \emph{ceiling} or \emph{interior} according to whether $x_j=l_j$, $x_j=h_j$ or
$l_j<x_j<h_j$. Let $\Pi$ be the set of $(x',f')\in\mathbb R^n\times\mathbb R^{m\times n}$ with
\begin{itemize}[nosep]
\item $x'_j=l_j$ on pinned and floor goods, $x'_j=h_j$ on ceiling goods, $l_j\le x'_j\le h_j$ on
interior goods;
\item $f'\ge0$, $f'_{ij}=0$ for $(i,j)\notin A$, and $\sum_jf'_{ij}=p_i$ for all $i$;
\item $a_{ij}-b_{ij}x'_j\ge a_{ik}-b_{ik}x'_k$ for all $(i,j)\in A$ and all $k$;
\item $\sum_if'_{ij}\le c_j$ on floor goods, $\sum_if'_{ij}\ge c_j$ on ceiling goods and
$\sum_if'_{ij}=c_j$ on interior goods.
\end{itemize}
\emph{$\Pi$ contains $(x,f)$.} The labels are read off $x$. A positive $f_{ij}$ has $(i,j)\in A$ and is a
best response, so the comparison rows hold. A floor good has $x_j=l_j<h_j$, so \eqref{eq:sxm-box} gives
$w_j\le c_j$; a ceiling good has $x_j=h_j>l_j$, so $w_j\ge c_j$; an interior good has both, so $w_j=c_j$.

\emph{Every point of $\Pi$ is an answer.} It lies in $[l,h]$ and has full rows. If $f'_{ij}>0$ then
$(i,j)\in A$, and the comparison rows say that $j$ maximizes $U_{ik}(x')$ over all $k$. For
\eqref{eq:sxm-box}: a pinned good has no condition; a floor good has $x'_j=l_j$, so only $x'_j<h_j$ can
apply, and then $\sum_if'_{ij}\le c_j$ holds; a ceiling good is symmetric; an interior good has load
exactly $c_j$, which satisfies both implications.

\emph{A vertex of $\Pi$ is short.} $\Pi$ is a nonempty polytope, since $x'\in[l,h]$ and
$0\le f'_{ij}\le p_i$. So it has a vertex, and one is computed by rational linear programming in time
polynomial in the length of the rows. Multiply each row by the product of the denominators of its
entries and right-hand side. A comparison row involves $b_{ij},b_{ik},a_{ij},a_{ik}$, so its integer entries
have absolute value below $2^{4\sigma+1}$; every other row involves one data number besides
coefficients $\pm1$, so its entries are below $2^{\sigma}$. A vertex is the unique solution of
$N_{\mathrm v}$ linearly independent rows of $\Pi$ satisfied with equality. By Cramer's rule each entry
is a quotient of two determinants of integer $N_{\mathrm v}\times N_{\mathrm v}$ matrices with entries
below $2^{4\sigma+1}$, and by Hadamard's inequality each determinant has absolute value at most
$(N_{\mathrm v}2^{4\sigma+1})^{N_{\mathrm v}}$. So numerator and denominator each have at most
$N_{\mathrm v}(\log_2N_{\mathrm v}+4\sigma+1)+1$ bits. Reading off $A$ and the labels from $(x,f)$ takes
time linear in its length.
\end{proof}

In particular every box query that has an answer has one of binary length polynomial in the query.
Every box query of a positive game has an answer (a BNE of $\widehat\Gamma$ exists by Nash's theorem,
and Lemmas~\ref{lem:sxm-game}(b), \ref{lem:sxm-bounded}(b) and~\ref{lem:sxm-reserve} below turn it into
an answer; alternatively, use Lemma~\ref{lem:mk-box-m1} together with the existence of (M1) equilibria,
Section~\ref{sec:pre-market}).

\paragraph{The reserve market of a query.} Fix a box query $(c,l,h)$ of a positive game and put
\[
v_i^0=\max_k\bigl(a_{ik}-b_{ik}h_k\bigr)\qquad(i\in[m]),
\]
as in Lemma~\ref{lem:mk-box-m1}. Let $\mathcal M$ be the dense (M1) market in general form
(Section~\ref{sec:pre-market}) with buyers $i\in[m]$ of mass $p_i$ and reserve $v_i^0$, goods $j\in[n]$
of capacity $c_j$ and reserve $l_j$, and payoff $a_{ij}-b_{ij}P_j$ to buyer $i$ from good $j$ at price
$P_j$. Written out, $(P,v,f)$ with $P\in\mathbb R^n$, $v\in\mathbb R^m$, $f\in\mathbb R^{m\times n}$ is an
equilibrium of $\mathcal M$ if and only if
\begin{enumerate}[nosep,label=(R\arabic*)]
\item $P_j\ge l_j$ for all $j$ and $v_i\ge v_i^0$ for all $i$;
\item $f\ge0$; $v_i\ge a_{ij}-b_{ij}P_j$ for all $(i,j)$, with equality if $f_{ij}>0$;
\item $\sum_jf_{ij}\le p_i$, with equality if $v_i>v_i^0$;
\item $\sum_if_{ij}\le c_j$, with equality if $P_j>l_j$.
\end{enumerate}
This is a matching market with singles in the sense of Galichon and Jacquet~\cite{galichon25}: buyer
mass left unmatched receives the reservation utility $v_i^0$, and capacity left unsold the reservation
price $l_j$.

\begin{lemma}[Singles]\label{lem:sxm-reserve}
Every equilibrium $(P,v,f)$ of $\mathcal M$ satisfies $l\le P\le h$. For each $i$ fix an index
$k(i)\in\argmax_k(a_{ik}-b_{ik}h_k)$, and put $x=P$ and
\[
f'_{ij}=f_{ij}+[j=k(i)]\Bigl(p_i-\sum_{k}f_{ik}\Bigr).
\]
Then $(x,f')$ is an answer to the query $(c,l,h)$.
\end{lemma}

\begin{proof}
This is Lemma~\ref{lem:mk-box-m1} before the shift to zero reserves; we repeat the argument in the
present notation. $P\ge l$ is (R1). Suppose $P_j>h_j$ for some $j$. For every $i$, since $b_{ij}>0$,
\[
a_{ij}-b_{ij}P_j<a_{ij}-b_{ij}h_j\le v_i^0\le v_i,
\]
so the equality in (R2) fails and $f_{ij}=0$. Column $j$ is then empty, while $P_j>h_j\ge l_j$ and
(R4) require $\sum_if_{ij}=c_j>0$. This contradiction gives $P\le h$.

Let $i$ have \emph{spare mass} $s_i=p_i-\sum_kf_{ik}>0$. By (R3), $v_i=v_i^0$. With $k=k(i)$, (R2)
gives $a_{ik}-b_{ik}h_k=v_i^0=v_i\ge a_{ik}-b_{ik}P_k$, so $P_k\ge h_k$, hence $P_k=h_k$, and the pair
$(i,k)$ satisfies (R2) with equality. So $f'$ is obtained from $f$ by adding mass only on pairs that are
tight in (R2), and only at goods with $x_k=h_k$.

Now check the answer conditions. $l\le x\le h$ was shown. $f'\ge0$ and $\sum_jf'_{ij}=p_i$ by
construction. If $f'_{ij}>0$, the pair $(i,j)$ is tight, so $a_{ij}-b_{ij}x_j=v_i\ge a_{ik}-b_{ik}x_k$
for all $k$ by (R2), and $j$ maximizes $U_{ik}(x)$. For \eqref{eq:sxm-box}: if $x_j>l_j$, then (R4) gives
$\sum_if_{ij}=c_j$, and $f'\ge f$ gives $w_j\ge c_j$; if $x_j<h_j$, no mass was added at $j$, so
$w_j=\sum_if_{ij}\le c_j$ by (R4).
\end{proof}

\begin{lemma}[Transposition]\label{lem:sxm-transpose}
Let $\mathcal M^\top$ be the dense (M1) market in general form with buyers $j\in[n]$ of mass $c_j$ and
reserve $l_j$, goods $i\in[m]$ of capacity $p_i$ and reserve $v_i^0$, and payoff
\[
\frac{a_{ij}}{b_{ij}}-\frac1{b_{ij}}\,t_i
\]
to buyer $j$ from a unit of good $i$ at price $t_i$ (slopes $1/b_{ij}>0$). Then $(P,v,f)$ is an
equilibrium of $\mathcal M$ if and only if $(v,P,f^\top)$ is an equilibrium of $\mathcal M^\top$, where $v$
is now the price vector, $P$ the utility vector, and $(f^\top)_{ji}=f_{ij}$ the flow.
\end{lemma}

\begin{proof}
Write the conditions of $\mathcal M^\top$ for $(t,u,g)=(v,P,f^\top)$ (prices $t$, utilities $u$, flows
$g$) and compare them with (R1)--(R4).
\begin{itemize}[nosep]
\item Reserves: $u_j\ge l_j$ and $t_i\ge v_i^0$, that is, $P_j\ge l_j$ and $v_i\ge v_i^0$: this is (R1).
\item No blocking: $u_j\ge(a_{ij}-t_i)/b_{ij}$, that is, $b_{ij}P_j\ge a_{ij}-v_i$ (multiply by
$b_{ij}>0$), that is, $v_i\ge a_{ij}-b_{ij}P_j$; with equality on the same pairs, and $g_{ji}>0$ if and
only if $f_{ij}>0$: this is (R2).
\item Buyer $j$'s row: $\sum_ig_{ji}=\sum_if_{ij}\le c_j$, with equality if $u_j>l_j$, that is, if
$P_j>l_j$: this is (R4).
\item Good $i$'s column: $\sum_jg_{ji}=\sum_jf_{ij}\le p_i$, with equality if $t_i>v_i^0$, that is, if
$v_i>v_i^0$: this is (R3).
\end{itemize}
Each condition of one market is equivalent to one condition of the other, so the two sets of
equilibria correspond exactly.
\end{proof}

\paragraph{Removing the singles.} Let $c_\Sigma=\sum_jc_j>0$ and define, for $i\in[m]$ and $j\in[n]$,
\begin{gather*}
\gamma_i=1+\max\Bigl(0,\ \max_{j}\bigl(a_{ij}-v_i^0-b_{ij}l_j\bigr)\Bigr),\qquad
\hat a_{ji}=\frac{a_{ij}-v_i^0}{b_{ij}}-l_j,\qquad \hat b_{ji}=\frac{\gamma_i}{b_{ij}}>0,\\
\mu_j=\frac{c_j}{2c_\Sigma},\qquad \hat c_i=\frac{p_i}{2c_\Sigma},\qquad \hat c_o=1 .
\end{gather*}
Since $v_i^0\ge a_{ij}-b_{ij}h_j$, we have $a_{ij}-v_i^0-b_{ij}l_j\le b_{ij}(h_j-l_j)$, so
$1\le\gamma_i\le1+\max_jb_{ij}$. Let $\mathcal B$ be the box market (Definition~\ref{def:mk-box}, read
with these buyers and goods) with buyers $j\in[n]$ of masses $\mu_j$, goods $[m]\cup\{o\}$ (the
\emph{real goods} $i\in[m]$ and one \emph{outside good} $o$) of capacities $\hat c$, price box
$[0,1]^{[m]\cup\{o\}}$, and payoffs
\[
\hat U_{ji}(\xi)=\hat a_{ji}-\hat b_{ji}\xi_i\quad(i\in[m]),\qquad \hat U_{jo}(\xi)=-\xi_o .
\]
So an equilibrium of $\mathcal B$ is a pair $(\xi,\varphi)$ with $\xi\in[0,1]^{[m]\cup\{o\}}$,
$\varphi_{js}\ge0$, $\sum_s\varphi_{js}=\mu_j$, $\varphi_{js}>0$ only if $s$ maximizes $\hat U_{js'}(\xi)$
over $s'\in[m]\cup\{o\}$, and, with $\hat w_s=\sum_j\varphi_{js}$,
$\xi_s>0\Rightarrow\hat w_s\ge\hat c_s$ and $\xi_s<1\Rightarrow\hat w_s\le\hat c_s$. In words: prices
and utilities are shifted so that both reserves become zero, masses and capacities are divided by
$2c_\Sigma$, the price of real good $i$ is measured in units of $\gamma_i$ above $v_i^0$, and the
buyers' reserves become one outside good that is never sold out. The point of $\gamma_i$ is the
identity
\begin{equation}\label{eq:sxm-shift}
\hat U_{ji}(\xi)=\frac{a_{ij}-t_i}{b_{ij}}-l_j\qquad\text{for }t_i=v_i^0+\gamma_i\xi_i ,
\end{equation}
which follows by expanding $\hat a_{ji}-\hat b_{ji}\xi_i$.

\begin{lemma}[Outside good and ceilings]\label{lem:sxm-bounded}
\begin{enumerate}[nosep,label=(\alph*)]
\item Every equilibrium $(\xi,\varphi)$ of $\mathcal B$ has $\xi_o=0$ and $\xi_i<1$ for all $i\in[m]$; hence
$\hat w_i\le\hat c_i$ for all $i\in[m]$, with equality if $\xi_i>0$.
\item \emph{Decoding.} For an equilibrium $(\xi,\varphi)$ of $\mathcal B$ put
\[
v_i=v_i^0+\gamma_i\xi_i,\qquad P_j=l_j+\max\Bigl(0,\ \max_{i\in[m]}\hat U_{ji}(\xi)\Bigr),\qquad
f_{ij}=2c_\Sigma\,\varphi_{ji}\qquad(i\in[m],\ j\in[n]).
\]
Then $(v,P,f^\top)$ is an equilibrium of $\mathcal M^\top$, and so $(P,v,f)$ is an equilibrium of
$\mathcal M$.
\item \emph{Encoding.} Conversely, let $(P,v,f)$ be an equilibrium of $\mathcal M$. Put
$\xi_i=(v_i-v_i^0)/\gamma_i$, $\xi_o=0$, $\varphi_{ji}=f_{ij}/(2c_\Sigma)$ and
$\varphi_{jo}=\mu_j-\sum_i\varphi_{ji}$. Then $(\xi,\varphi)$ is an equilibrium of $\mathcal B$, and
decoding it by (b) returns $(P,v,f)$.
\end{enumerate}
\end{lemma}

\paragraph{Proof roadmap.} (a) The outside good cannot fill its capacity, so its price is zero; a real
good at its ceiling is worse than the outside good for every buyer, so it would be unsold, which a
positive price forbids. (b) and (c) are the conditions of $\mathcal M^\top$ read through
\eqref{eq:sxm-shift}, in both directions.

\begin{proof}
(a) All buyers together have mass $\sum_j\mu_j=1/2$, so $\hat w_o\le1/2<1=\hat c_o$, and $\xi_o>0$ would
require $\hat w_o\ge\hat c_o$. So $\xi_o=0$, $\hat U_{jo}=0$, and buyer $j$'s best value is
\[
\hat V_j=\max\Bigl(0,\ \max_{i\in[m]}\hat U_{ji}(\xi)\Bigr)\ \ge0 .
\]
Suppose $\xi_i=1$ for some $i\in[m]$. By \eqref{eq:sxm-shift} with $t_i=v_i^0+\gamma_i$, and by
$\gamma_i>a_{ij}-v_i^0-b_{ij}l_j$ and $b_{ij}>0$,
\[
\hat U_{ji}(\xi)=\frac{a_{ij}-v_i^0-\gamma_i-b_{ij}l_j}{b_{ij}}<0\le\hat V_j\qquad\text{for every }j .
\]
So no buyer uses $i$ and $\hat w_i=0<\hat c_i$, as $p_i>0$; but $\xi_i=1>0$ requires
$\hat w_i\ge\hat c_i$. Hence $\xi_i<1$, so $\hat w_i\le\hat c_i$, and if $\xi_i>0$ also
$\hat w_i\ge\hat c_i$.

(b) Let $t_i=v_i$, so that \eqref{eq:sxm-shift} reads $\hat U_{ji}(\xi)=(a_{ij}-v_i)/b_{ij}-l_j$, and
$P_j-l_j=\hat V_j$ by (a). We verify the conditions of $\mathcal M^\top$ for prices $v$, utilities $P$
and flows $g_{ji}=f_{ij}$.
\begin{itemize}[nosep]
\item Reserves: $P_j\ge l_j$ since $\hat V_j\ge0$; $v_i\ge v_i^0$ since $\xi_i\ge0$.
\item No blocking: $P_j-l_j=\hat V_j\ge\hat U_{ji}(\xi)$, that is, $P_j\ge(a_{ij}-v_i)/b_{ij}$.
\item Tightness: $g_{ji}>0$ gives $\varphi_{ji}>0$, so $\hat U_{ji}(\xi)=\hat V_j$, that is,
$P_j=(a_{ij}-v_i)/b_{ij}$.
\item Buyer $j$'s row: $\sum_ig_{ji}=2c_\Sigma(\mu_j-\varphi_{jo})=c_j-2c_\Sigma\varphi_{jo}\le c_j$. If
$P_j>l_j$, then $\hat V_j>0=\hat U_{jo}$, so $\varphi_{jo}=0$ and the row is full.
\item Good $i$'s column: $\sum_jg_{ji}=2c_\Sigma\hat w_i\le2c_\Sigma\hat c_i=p_i$ by (a). If
$v_i>v_i^0$, then $\xi_i>0$, so $\hat w_i=\hat c_i$ by (a) and the column is full.
\end{itemize}
Lemma~\ref{lem:sxm-transpose} turns this into an equilibrium $(P,v,f)$ of $\mathcal M$.

(c) By Lemma~\ref{lem:sxm-transpose}, $(v,P,f^\top)$ is an equilibrium of $\mathcal M^\top$. Two facts
about it come first.

\emph{The utility is a maximum:} $P_j=\max\bigl(l_j,\max_i(a_{ij}-v_i)/b_{ij}\bigr)$. The inequality
$\ge$ is (R1) with (R2). If $P_j>l_j$, then (R4) fills the row of $j$ in $\mathcal M^\top$; since
$c_j>0$ some $f_{ij}>0$, and that pair is tight.

\emph{The price is below the ceiling:} $v_i-v_i^0<\gamma_i$. This is clear if $v_i=v_i^0$. If
$v_i>v_i^0$, (R3) fills the column of $i$; since $p_i>0$ some $f_{ij}>0$, and tightness with $P_j\ge l_j$
gives $(a_{ij}-v_i)/b_{ij}=P_j\ge l_j$, that is,
$v_i-v_i^0\le a_{ij}-v_i^0-b_{ij}l_j<\gamma_i$.

So $\xi\in[0,1)^m\times\{0\}$. The flows are nonnegative: $\varphi_{ji}\ge0$, and
$\varphi_{jo}=(c_j-\sum_if_{ij})/(2c_\Sigma)\ge0$ by (R4); the rows sum to $\mu_j$ by definition. By
\eqref{eq:sxm-shift} with $t_i=v_i$, $\hat U_{ji}(\xi)=(a_{ij}-v_i)/b_{ij}-l_j$, and $\hat U_{jo}=0$; so
by the first fact the best value of $j$ is $P_j-l_j$. A real good used by $j$ has $f_{ij}>0$, hence is
tight and attains $P_j-l_j$. If $\varphi_{jo}>0$, the row of $j$ is not full, so $P_j=l_j$ by (R4), and
$o$ attains the best value $0$. Loads: $\hat w_i=\sum_jf_{ij}/(2c_\Sigma)\le p_i/(2c_\Sigma)=\hat c_i$ by
(R3), with equality if $\xi_i>0$, because then $v_i>v_i^0$; since $\xi_i<1$ both box conditions hold. For
$o$, $\xi_o=0<1$ and $\hat w_o\le1/2\le\hat c_o$. Finally, decoding returns $v_i^0+\gamma_i\xi_i=v_i$,
$2c_\Sigma\varphi_{ji}=f_{ij}$, and $l_j+\max(0,\max_i\hat U_{ji}(\xi))=\max(l_j,\max_i(a_{ij}-v_i)/b_{ij})=P_j$.
\end{proof}

\paragraph{The game of a query.} Put $K=m+3$. The game $\widehat\Gamma$ has the following data
(Definition~\ref{def:game}). In it the index $j$ runs over types and $i$ over targets.
\begin{itemize}[nosep]
\item \emph{Targets:} the real targets $i\in[m]$, the outside target $o$, and auxiliary targets
$\mathrm{aux}_1,\dots,\mathrm{aux}_K$; in total $2m+4$ targets.
\item \emph{Resources:} $\hat r=m+2$.
\item \emph{Types:} real types $j\in[n]$ with prior $\mu_j=c_j/(2c_\Sigma)$, and dedicated types
$\mathrm{ded}_1,\dots,\mathrm{ded}_K$ with prior $1/(2K)$ each; in total $n+m+3$ types, with positive
priors summing to $1$.
\item \emph{Defender:} payoff $0$ at an uncovered target, and gains $\hat d_i=1/\hat c_i=2c_\Sigma/p_i$
at real targets, $\hat d_o=1/\hat c_o=1$, and $\hat d_{\mathrm{aux}_k}=2K$.
\item \emph{Attackers:} with coverage $\xi$, real type $j$ receives $\hat a_{ji}-\hat b_{ji}\xi_i$ at real
target $i$, $-\xi_o$ at $o$, and $-2-\xi_{\mathrm{aux}_k}$ at $\mathrm{aux}_k$; dedicated type
$\mathrm{ded}_k$ receives $2-\xi_{\mathrm{aux}_k}$ at $\mathrm{aux}_k$ and $-2-\xi_s$ at every other target
$s$.
\end{itemize}
Write $\varphi_{\theta s}$ for the attack mass of type $\theta$ on target $s$, $\hat w_s$ for the load,
and $\hat\pi_s=\hat d_s\hat w_s$ for the score.

\begin{lemma}[The auxiliary targets force the cutoff]\label{lem:sxm-game}
\begin{enumerate}[nosep,label=(\alph*)]
\item $\widehat\Gamma$ is a Bayesian simple security game with $2m+4$ targets, $n+m+3$ types and
$\hat r=m+2$ resources, $0<\hat r<2m+4$, all attacker slopes positive. If $\sigma\ge1$ bounds the
binary length of every number of the query and the game, every number of $\widehat\Gamma$ has binary
length $O(n\sigma)$, and $\widehat\Gamma$ is computed in time polynomial in $n$, $m$ and $\sigma$.
\item Let $(\xi,\varphi)$ be a BNE of $\widehat\Gamma$ (coverage and attack masses). Then every dedicated
type $\mathrm{ded}_k$ attacks only $\mathrm{aux}_k$, no real type attacks an auxiliary target, every
auxiliary score equals $1$, $1\le\sum_k\xi_{\mathrm{aux}_k}\le m+2<K$, and every cutoff $\hat\lambda$ of
Lemma~\ref{lem:cutoff} for this BNE equals $1$. At a real target and at $o$ the score is the load divided
by the capacity, $\hat\pi_s=\hat w_s/\hat c_s$, and the restriction of $(\xi,\varphi)$ to the goods
$[m]\cup\{o\}$ and the real types is an equilibrium of $\mathcal B$.
\item Conversely, every equilibrium of $\mathcal B$ extends to a BNE of $\widehat\Gamma$.
\end{enumerate}
\end{lemma}

\paragraph{Proof roadmap.} (b) has four steps: (1) payoffs separate the type classes, so auxiliary loads
are constant and auxiliary scores are $1$; (2) counting the budget puts the auxiliary coverage strictly
inside $(0,K)$; (3) one auxiliary target with positive coverage and one below full coverage pin the
cutoff to $1$; (4) at cutoff $1$ the defender's conditions on the other targets are the box conditions
of $\mathcal B$.

\begin{proof}
(a) All attacker slopes ($\hat b_{ji}$ and $1$) are positive, so $u_2^0\ge u_2^1$ holds at every pair;
all gains are positive, so $u_1^0=0\le u_1^1$. The priors are positive with
$\sum_j\mu_j+K/(2K)=1/2+1/2=1$, and $0<m+2<2m+4$. For the lengths, $v_i^0$ and the inner maximum in
$\gamma_i$ are each the maximum of $n$ numbers, each obtained from at most six query numbers by two
products and three additions, so they and $\gamma_i$, $\hat a_{ji}$, $\hat b_{ji}$ have binary length
$O(\sigma)$; the sum $c_\Sigma$ of $n$ rationals has binary length $O(n\sigma)$, hence so do $\mu_j$ and
$\hat d_i$; the remaining numbers are integers or $1/(2K)$, of length $O(\log m)\subseteq O(\sigma)$
because some $p_i\le1/m$. The construction uses $O((m+n)m)$ arithmetic operations.

(b) \emph{Separation.} A dedicated type $\mathrm{ded}_k$ gets $2-\xi_{\mathrm{aux}_k}\ge1$ at
$\mathrm{aux}_k$ and at most $-2$ elsewhere, so it attacks only $\mathrm{aux}_k$. A real type gets
$-\xi_o\ge-1$ at $o$ and at most $-2$ at every auxiliary target, so it never attacks an auxiliary target,
and its best value is attained on $[m]\cup\{o\}$. Hence $\hat w_{\mathrm{aux}_k}=1/(2K)$ and
$\hat\pi_{\mathrm{aux}_k}=2K/(2K)=1$, and the loads of the real targets and of $o$ come from real types
only.

\emph{Counting.} The $m+1$ targets in $[m]\cup\{o\}$ carry coverage at most $m+1$ in total, and the total
coverage is $\hat r=m+2$. So $\sum_k\xi_{\mathrm{aux}_k}\ge1>0$, and also
$\sum_k\xi_{\mathrm{aux}_k}\le m+2<K$. Hence some auxiliary target has $\xi_{\mathrm{aux}_k}>0$ and some
has $\xi_{\mathrm{aux}_{k'}}<1$.

\emph{Cutoff.} Since $0<\hat r<2m+4$, Lemma~\ref{lem:cutoff} gives a cutoff $\hat\lambda$ with
$\xi_s>0\Rightarrow\hat\pi_s\ge\hat\lambda$ and $\xi_s<1\Rightarrow\hat\pi_s\le\hat\lambda$. Applied to
$\mathrm{aux}_k$ and $\mathrm{aux}_{k'}$ this gives $1\ge\hat\lambda$ and $1\le\hat\lambda$. So every such
$\hat\lambda$ equals $1$.

\emph{Box conditions.} For $s\in[m]\cup\{o\}$, $\hat d_s=1/\hat c_s$ gives $\hat\pi_s=\hat w_s/\hat c_s$.
With $\hat\lambda=1$ the cutoff conditions read $\xi_s>0\Rightarrow\hat w_s\ge\hat c_s$ and
$\xi_s<1\Rightarrow\hat w_s\le\hat c_s$, which are the box conditions of $\mathcal B$ on $[0,1]$. A real
type $j$ has total mass $\mu_j$, all of it on $[m]\cup\{o\}$, on targets that maximize its payoff over all
targets and hence over $[m]\cup\{o\}$, where its payoffs are those of $\mathcal B$. So the restriction is
an equilibrium of $\mathcal B$.

(c) Let $(\xi,\varphi)$ be an equilibrium of $\mathcal B$. By Lemma~\ref{lem:sxm-bounded}(a),
$\xi_o=0$ and $\xi_i<1$, so $\Sigma:=\sum_{s\in[m]\cup\{o\}}\xi_s<m$. Give every auxiliary target the
coverage $\zeta=(m+2-\Sigma)/K$, which lies in $(2/K,(m+2)/K]\subseteq(0,1)$, so the total coverage is
$m+2$ and every coverage lies in $[0,1]$; let $\mathrm{ded}_k$ attack $\mathrm{aux}_k$ and the real types
play $\varphi$. Attackers best respond: the dedicated types as in (b), and the real types because their
payoffs at auxiliary targets are at most $-2<0=-\xi_o$, so their best value over all targets is the one
in $\mathcal B$. For the defender take the cutoff $\hat\lambda=1$: auxiliary scores equal $1$, and on
$[m]\cup\{o\}$ the cutoff conditions are the box conditions of $\mathcal B$, as in (b). By
Lemma~\ref{lem:cutoff}(b) and Lemma~\ref{lem:realize} this is a BNE.
\end{proof}

\begin{proposition}[$L^{O(\sqrt{\min(m,n)})}$]\label{prop:sx-min}
\statusrefereed
An exact BNE of an explicitly given Bayesian simple security game with $m$ types, $n$ targets and binary
input length $L\ge2$ can be computed deterministically in $L^{O(\sqrt{\min(m,n)})}$ bit operations. Zero
attacker slopes $b_{ij}$ and zero defender gains $d_j$ are allowed. Since $mn\le L$, this is
$2^{O(L^{1/4}\log L)}$ bit operations. For $m<n$ the algorithm is the game-to-market reduction of
Section~\ref{sec:mk-g2m} (Proposition~\ref{prop:mk-g2m}), whose box queries
(Lemma~\ref{lem:sxm-interface}) are answered through games with $2m+4$ targets
(Lemma~\ref{lem:sxm-game}) solved by Theorem~\ref{thm:sx-main}.
\end{proposition}

The bound depends on the separate game-to-market reduction of Chapter~\ref{ch:market}, through
Lemma~\ref{lem:sxm-interface}. Theorem~\ref{thm:sx-main} enters twice: directly when $n\le m$, and as the
solver of the query games when $m<n$.

\paragraph{Proof roadmap.} (1) Delete zero-prior types and settle $r\in\{0,n\}$ and $n\le m$ directly.
(2) Make slopes and gains positive by the perturbation $\epsilon=g/4$ and repair at the end with one
pattern LP of Section~\ref{sec:contraction}. (3) For a positive game with $m<n$, run the procedure of
Lemma~\ref{lem:sxm-interface}, and answer each query through $\mathcal M$, $\mathcal M^\top$,
$\mathcal B$ and $\widehat\Gamma$ (Lemmas~\ref{lem:sxm-reserve}--\ref{lem:sxm-game}), solving
$\widehat\Gamma$ by Theorem~\ref{thm:sx-main} and shortening the answer by Lemma~\ref{lem:sxm-vertex}.
(4) Count: polynomially many queries, each in $L^{O(\sqrt m)}$, plus polynomial work.

\begin{proof}
\emph{Step 1: reductions.} Delete the types of prior zero; they are given any best response at the end
(Section~\ref{sec:pre-game}). Let $m'\le m$ be the number of remaining types; their priors still sum to
$1$. If $r\in\{0,n\}$ the coverage is forced and every type attacks a best target. If $n\le m'$, apply
Theorem~\ref{thm:sx-main}, which allows zero slopes and gains, in $L^{O(\sqrt n)}$ bit operations; here
$n=\min(m,n)$. From now on $0<r<n$, all priors are positive, and $m'<n$, so $m'\le\min(m,n)$. To keep the
notation, write $m$ for $m'$.

\emph{Step 2: zero slopes or gains.} If every $b_{ij}$ and every $d_j$ is positive, go to Step~3 with
the game itself. Otherwise let $Q$ be the constant of Lemma~\ref{lem:ueopl-pattern}(b) for the game,
computable in polynomial time with $\log Q$ polynomial in $L$, put $g=1/(4Q)$ and $\epsilon=g/4$, and let
$\Gamma_\epsilon$ be the game with $b^\epsilon_{ij}=b_{ij}+\epsilon$ and $d^\epsilon_j=d_j+\epsilon$ and the
same $p$, $a$ and $r$. It is a positive game with the same $m$ and $n$, and its binary length $L_\epsilon$
is polynomial in $L$. Let $(x,f)$ be an exact BNE of $\Gamma_\epsilon$ (coverage and attack masses),
computed by Step~3. By Lemma~\ref{lem:cutoff}(a) it has a cutoff
$\lambda^\epsilon\in[0,\max_jd^\epsilon_j]$. Take the pattern of Section~\ref{sec:contraction} with
allowed edges $\{(i,j):f_{ij}>0\}$ and labels zero, one or free according to $x_j=0$, $x_j=1$ or
$0<x_j<1$, and the point $(x,f,v,\lambda^\epsilon)$ with $v_i=\max_kU_{ik}(x)$ for the \emph{original}
payoffs $U_{ik}(x)=a_{ik}-b_{ik}x_k$. We bound its residual in every constraint of this pattern LP of the
original game.
\begin{itemize}[nosep]
\item Exactly satisfied: the budget, $0\le x\le1$, the row sums, $f_{ij}=0$ off the allowed edges,
$v_i\ge U_{ij}(x)$ for all pairs, $x_j=0$ at zero targets and $x_j=1$ at one targets.
\item Best responses on allowed edges: $f_{ij}>0$ means
$a_{ij}-b^\epsilon_{ij}x_j=\max_k(a_{ik}-b^\epsilon_{ik}x_k)$. Since
$U_{ik}(x)=a_{ik}-b^\epsilon_{ik}x_k+\epsilon x_k$ and $0\le x\le1$,
$v_i\le\max_k(a_{ik}-b^\epsilon_{ik}x_k)+\epsilon=a_{ij}-b^\epsilon_{ij}x_j+\epsilon\le U_{ij}(x)+\epsilon$. So
$0\le v_i-U_{ij}(x)\le\epsilon$.
\item Cutoff conditions: $d_jw_j=d^\epsilon_jw_j-\epsilon w_j$ with $0\le w_j\le1$. At a zero target
$x_j<1$, so $d_jw_j\le d^\epsilon_jw_j\le\lambda^\epsilon$ exactly. At a one target $x_j>0$, so
$d_jw_j\ge\lambda^\epsilon-\epsilon$. At a free target both hold, so
$\lambda^\epsilon-\epsilon\le d_jw_j\le\lambda^\epsilon$.
\item Outer boxes, exactly: $x\in[0,1]^n$; $0\le f_{ij}\le p_i\le1$;
$0\le\lambda^\epsilon\le\max_jd_j+\epsilon\le1+\max_jd_j$ since $\epsilon\le1$; and
$|v_i|\le\max_j(|a_{ij}|+b_{ij})$.
\end{itemize}
Every residual is at most $\epsilon=g/4<g$. By Lemma~\ref{lem:ueopl-pattern}(c), in the form with the
slack restricted to $[0,1]$ noted after its proof, the unrelaxed pattern LP of the original game is
feasible, and by Lemma~\ref{lem:ueopl-pattern}(a) any rational feasible point, computed by one LP in
polynomial time, is an exact BNE of the original game. In either case, realize the coverage of the BNE
of the original game by Lemma~\ref{lem:realize}, and give the deleted types of prior zero any best
response.

\emph{Step 3: a positive game with $m<n$.} Run the procedure of Lemma~\ref{lem:sxm-interface} and answer
each query $(c,l,h)$ as follows.
\begin{enumerate}[nosep,label=(\roman*)]
\item Build $\widehat\Gamma$ from the query (Lemma~\ref{lem:sxm-game}(a)).
\item Compute an exact BNE of $\widehat\Gamma$ by Theorem~\ref{thm:sx-main}, and from it the coverage
vector $\xi$ and the attack masses $\varphi_{\theta s}$ (prior times conditional probability).
\item By Lemma~\ref{lem:sxm-game}(b) the restriction of $(\xi,\varphi)$ is an equilibrium of
$\mathcal B$; decode it to an equilibrium $(P,v,f)$ of $\mathcal M$ by
Lemma~\ref{lem:sxm-bounded}(b), and turn it into an answer $(x,f')$ by Lemma~\ref{lem:sxm-reserve}.
\item Replace $(x,f')$ by the short answer of Lemma~\ref{lem:sxm-vertex}, and return it.
\end{enumerate}
Each returned pair is an answer to its query, so by Lemma~\ref{lem:sxm-interface}(c) the procedure stops
with an exact BNE of the positive game. Every step is deterministic.

\emph{Step 4: running time.} Let $L_0$ be the binary length of the positive game of Step~3 (the game
itself or $\Gamma_\epsilon$); then $2\le L_0\le L^{c_0}$ for a constant $c_0$. By
Lemma~\ref{lem:sxm-interface}(b) there are at most $L_0^{c_1}$ queries, and each query, with the game
data, has binary length at most $L_0^{c_1}$, for a constant $c_1$. By Lemma~\ref{lem:sxm-game}(a), each
$\widehat\Gamma$ has binary length $\hat L$ with $2\le\hat L\le L_0^{c_2}$ and $2m+4\le6m$ targets.
Theorem~\ref{thm:sx-main} gives a constant $C$ such that (ii) takes at most
$\hat L^{C\sqrt{2m+4}}\le L_0^{\sqrt6\,c_2C\sqrt m}$ bit operations; this also bounds the binary length of
the BNE it outputs. Steps (i), (iii) and (iv) take time polynomial in $L_0$ and in that length
(Lemmas~\ref{lem:sxm-game}(a), \ref{lem:sxm-bounded}, \ref{lem:sxm-reserve} and~\ref{lem:sxm-vertex}).
So each query costs $L_0^{O(\sqrt m)}$ bit operations. Every returned answer has binary length polynomial
in $L_0$ by Lemma~\ref{lem:sxm-vertex}, so by Lemma~\ref{lem:sxm-interface}(d) the rest of the
procedure takes time polynomial in $L_0$. Step~2 adds the computation of $Q$ and one LP, Step~1 and the
realization add polynomial time. In total, with $L_0\le L^{c_0}$ and $m\le\min(m,n)$ (for the original
$m$ and $n$),
\[
L_0^{c_1}\cdot L_0^{O(\sqrt m)}+L^{O(1)}=L^{O(\sqrt{\min(m,n)})} .
\]

\emph{Step 5: the bound in $L$.} The input lists the $2mn$ attacker utilities $u_2^0(i,j)$ and
$u_2^1(i,j)$, each of binary length at least one, so $mn\le L$ and
$\min(m,n)\le\sqrt{mn}\le\sqrt L$. Hence
$L^{O(\sqrt{\min(m,n)})}\le L^{O(L^{1/4})}=2^{O(L^{1/4}\log L)}$.
\end{proof}

The constants of $\widehat\Gamma$ are not tight: the proof uses $K=m+3$ and $\hat r=m+2$ only through
$1\le\sum_k\xi_{\mathrm{aux}_k}<K$, and $\hat c_o=1$ only through $\hat c_o>\sum_j\mu_j$.
